% chapter22 —— 由 tools/convert.py 自动生成，勿手改
% 源文件：part0000_split_177.html, part0000_split_178.html, part0000_split_179.html, part0000_split_180.html, part0000_split_181.html, part0000_split_182.html, part0000_split_183.html, part0000_split_184.html, part0000_split_185.html, part0000_split_186.html, part0000_split_187.html, part0000_split_188.html

\begin{myEpigraph}{——《美狄亚》，Euripides}
希望没有人认为我卑微、软弱或顺从，

希望他们明白我与众不同：

对于敌人，我意味着危险；对于朋友，我意味着忠诚。

这便是我荣耀的人生。
\end{myEpigraph}

系统编程的隐秘乐趣在于，在每一种安全语言和精心设计的抽象之下，都是极度不安全的机器语言和按位操作的汹涌暗流。你也可以用 Rust 写出这种代码。

迄今为止，通过类型检查、生命周期检查、限界检查等方法，本书中介绍的这门语言可以确保你的程序完全自动地摆脱内存错误和数据竞争的困扰。但是这种自动化推理有其局限性，因为 Rust 中仍然有许多无法识别为安全的高价值技术。

\textbf{不安全}\footnote{此处的不安全（unsafe）并不表示一定有风险，只是说 Rust 编译器无法保证安全。unsafe 的程度相比中文的“不安全”要弱得多。但这是社区普遍采纳的一个翻译，因此本书跟随这种译法。——译者注}\textbf{代码}能让你告诉 Rust：“我选择使用你无法保证安全的特性。”通过将块或函数标记为不安全的，你可以获得调用标准库中的 \texttt{unsafe} 函数、解引用不安全指针以及调用以其他语言（如 C 和 C++）编写的函数等能力。Rust 的其他安全检查仍然适用：类型检查、生命周期检查和索引的边界检查都会正常进行。不安全代码只会启用一小部分附加特性。

这种跨越 Rust 安全边界的能力使得 Rust 可以实现自身许多最基本的特性，就像 C 和 C++ 被用于实现自己的标准库一样。使用不安全代码，\texttt{Vec} 类型可以更加高效地管理其缓冲区，\texttt{std::io} 模块可以和操作系统对话，\texttt{std::thread} 模块和 \texttt{std::sync} 模块可以提供并发原语。

本章涵盖了使用不安全特性的所有要点。

\begin{itemize}
\item Rust 的 \texttt{unsafe} 块在普通的、安全的 Rust 代码和使用了不安全特性的代码之间建立了边界。
\item 可以将函数标记为 \texttt{unsafe}，提醒调用者这里存在必须遵守的额外契约，以避免未定义行为。
\item 裸指针及其方法允许不受限制地访问内存，进而构建 Rust 的类型系统原本会禁止的数据结构。Rust 的引用是安全但受限的，而任何 C 或 C++ 程序员都知道，裸指针是一个强大而锋利的工具。
\item 理解未定义行为的定义将帮助你理解为什么它会产生比得到不正确的结果还要严重的后果。
\item 不安全特型（\texttt{unsafe trait}）与 \texttt{unsafe} 函数类似，对每个实现而不是每个调用者都强加了必须遵守的契约。
\end{itemize}

\section{不安全因素来自哪里}

在本书的开头，我们展示过一个因为没有遵守 C 标准规定中的规则而以令人惊讶的方式崩溃的 C 程序。在 Rust 中可以做到同样的事情：

\begin{Verbatim}[commandchars=\\\{\},breaklines=true,breakanywhere=true,fontsize=\codefont,xleftmargin=17bp]
\PYGZdl{}\PYG{+w}{ }cat\PYG{+w}{ }crash.rs
fn\PYG{+w}{ }main\PYG{o}{(}\PYG{o}{)}\PYG{+w}{ }\PYG{o}{\PYGZob{}}
\PYG{+w}{    }\PYG{n+nb}{let}\PYG{+w}{ }mut\PYG{+w}{ }a:\PYG{+w}{ }\PYG{n+nv}{usize}\PYG{+w}{ }\PYG{o}{=}\PYG{+w}{ }\PYG{l+m}{0}\PYG{p}{;}
\PYG{+w}{    }\PYG{n+nb}{let}\PYG{+w}{ }\PYG{n+nv}{ptr}\PYG{+w}{ }\PYG{o}{=}\PYG{+w}{ }\PYG{p}{\PYGZam{}}mut\PYG{+w}{ }a\PYG{+w}{ }as\PYG{+w}{ }*mut\PYG{+w}{ }usize\PYG{p}{;}
\PYG{+w}{    }unsafe\PYG{+w}{ }\PYG{o}{\PYGZob{}}
\PYG{+w}{        }*ptr.offset\PYG{o}{(}\PYG{l+m}{3}\PYG{o}{)}\PYG{+w}{ }\PYG{o}{=}\PYG{+w}{ }0x7ffff72f484c\PYG{p}{;}
\PYG{+w}{    }\PYG{o}{\PYGZcb{}}
\PYG{o}{\PYGZcb{}}
\PYGZdl{}\PYG{+w}{ }cargo\PYG{+w}{ }build
\PYG{+w}{   }Compiling\PYG{+w}{ }unsafe\PYGZhy{}samples\PYG{+w}{ }v0.1.0
\PYG{+w}{    }Finished\PYG{+w}{ }debug\PYG{+w}{ }\PYG{o}{[}unoptimized\PYG{+w}{ }+\PYG{+w}{ }debuginfo\PYG{o}{]}\PYG{+w}{ }target\PYG{o}{(}s\PYG{o}{)}\PYG{+w}{ }\PYG{k}{in}\PYG{+w}{ }\PYG{l+m}{0}.44s
\PYGZdl{}\PYG{+w}{ }../../target/debug/crash
crash:\PYG{+w}{ }Error:\PYG{+w}{ }.netrc\PYG{+w}{ }file\PYG{+w}{ }is\PYG{+w}{ }readable\PYG{+w}{ }by\PYG{+w}{ }others.
crash:\PYG{+w}{ }Remove\PYG{+w}{ }password\PYG{+w}{ }or\PYG{+w}{ }make\PYG{+w}{ }file\PYG{+w}{ }unreadable\PYG{+w}{ }by\PYG{+w}{ }others.
Segmentation\PYG{+w}{ }fault\PYG{+w}{ }\PYG{o}{(}core\PYG{+w}{ }dumped\PYG{o}{)}
\PYGZdl{}
\end{Verbatim}

这个程序借用了对局部变量 \texttt{a} 的可变引用，将其转换为 \texttt{*mut usize} 类型的裸指针，然后使用 \texttt{offset} 方法在内存中又生成了 3 个字的指针。这恰好是存储 \texttt{main} 的返回地址的地方。这个程序用一个常量覆盖了返回地址，这样从 \texttt{main} 返回的行为就会令人非常惊讶。导致这次崩溃的原因是程序错误地使用了不安全特性——在这个例子中就是解引用裸指针的能力。

不安全特性是强加了某种\textbf{契约}的特性：Rust 不能自动执行这些规则，但你必须遵守这些规则以避免\textbf{未定义行为}。

这种契约超出了常规类型检查和生命周期检查的能力范围，针对该不安全特性强加了更多规则。通常，Rust 本身根本不了解契约，契约只是在该特性的文档中进行了解释。例如，裸指针类型有一个契约，它禁止解引用已超出其原始引用目标末尾的指针。上述例子中的表达式 \texttt{*ptr.offset(3) = ...} 破坏了这个契约。但是，正如前面的记录所示，Rust 毫无怨言地编译了这段程序，因为它的安全检查并未检测到这种违规行为。当使用了不安全特性时，作为程序员，你有责任检查自己的代码是否遵守了它们的契约。

许多特性需要遵守某些规则才能正确使用，但这些规则并不是这里所说的契约，除非违反它们的后果包括未定义行为。未定义行为是“Rust 坚定地认为你的代码永远不会出现的行为”。例如，Rust 认为你不会用其他内容覆盖函数调用的返回地址。能够通过 Rust 通常的安全检查并遵守其用到的不安全特性的契约的代码不可能做这样的事情。由于前面的程序违反了裸指针契约，因此其行为是未定义的，它已经偏离了轨道。

如果代码表现出未定义行为，那你就已经违背了与 Rust 达成的交易，所以 Rust 无法对其后果负责。从系统库深处挖掘出不相关的错误消息并导致崩溃是一种可能的后果，将计算机的控制权交给攻击者是另一种后果。在没有警告的情况下，从 Rust 的一个版本换到下一个版本可能会产生不同的效果。然而，有时未定义行为并没有明显的后果。如果 \texttt{main} 函数永远不会返回（比如调用了 \texttt{std::process::exit} 来提前终止程序），那么损坏的返回地址可能无关紧要。

只能在 \texttt{unsafe} 块或 \texttt{unsafe} 函数中使用不安全特性，我们将在接下来的内容中对两者进行解释。这可以避免在不知不觉中使用不安全特性：通过强制编写 \texttt{unsafe} 块或函数，Rust 会确保你已经知道在自己的代码中可能要遵守的额外规则。

\section{不安全块}

\texttt{unsafe} 块看起来就像前面加了 \texttt{unsafe} 关键字的普通 Rust 块，不同之处在于可以在块中使用不安全特性：

\begin{Verbatim}[commandchars=\\\{\},breaklines=true,breakanywhere=true,fontsize=\codefont,xleftmargin=17bp]
\PYG{k}{unsafe}\PYG{+w}{ }\PYG{p}{\PYGZob{}}
\PYG{+w}{    }\PYG{n+nb}{String}\PYG{p}{::}\PYG{n}{from\PYGZus{}utf8\PYGZus{}unchecked}\PYG{p}{(}\PYG{n}{ascii}\PYG{p}{)}
\PYG{p}{\PYGZcb{}}
\end{Verbatim}

如果块前面没有 \texttt{unsafe} 关键字，那么 Rust 就会反对使用 \texttt{from\_utf8\_unchecked}，因为这是一个 \texttt{unsafe} 函数。有了它周围的 \texttt{unsafe} 块，就可以在任何地方使用此代码了。

与普通的 Rust 块一样，\texttt{unsafe} 块的值就是其最终表达式的值，如果没有则为 \texttt{()}。前面展示的对 \texttt{String::from\_utf8\_unchecked} 的调用提供了该块的值。

\texttt{unsafe} 块解锁了 5 个额外的选项。

\begin{itemize}
\item 可以调用 \texttt{unsafe} 函数。每个 \texttt{unsafe} 函数都必须根据自己的目的指定自己的契约。
\item 可以解引用裸指针。安全代码可以传递裸指针，比较它们，并从引用（甚至整数）转换成它们，但只有不安全代码才能真正使用它们来访问内存。22.8 节将详细介绍裸指针并解释如何安全地使用它们。
\item 可以访问 \texttt{union} 的各个字段，编译器无法确定这些字段是否包含其各自类型的有效位模式。
\item 可以访问可变的 \texttt{static} 变量。如 19.3.11 节所述，Rust 无法确定线程何时使用可变 \texttt{static} 变量，因此它们的契约要求你确保所有访问都能正确同步。
\item 可以访问通过 Rust 的外部函数接口声明的函数和变量。即使声明为不可变的，这些函数和变量也仍然会被看作 \texttt{unsafe} 的，因为它们对于用其他可能不遵守 Rust 安全规则的语言编写的代码仍然是可见的。
\end{itemize}

将不安全特性限制在 \texttt{unsafe} 块中并不能真正阻止你做任何想做的事。你完全可以只将一个 \texttt{unsafe} 块粘贴到代码中，然后继续我行我素。该规则的主要目的在于将人们的视线吸引到 Rust 无法保证其安全性的代码上。

\begin{itemize}
\item 你不会无意中使用不安全特性，然后发现要对连自己都不知道在哪里的契约负责。
\item \texttt{unsafe} 块会引起评审者的更多关注。有些项目甚至会通过自动化设施来确保这一点，它们会标记出影响 \texttt{unsafe} 块的代码更改以引起特别关注。
\item 当你考虑编写 \texttt{unsafe} 块时，可以花点儿时间问问自己是否真的需要这样的措施。如果是为了性能，那是否有测量结果表明这确实是一个瓶颈呢？也许在安全的 Rust 中有更好的办法来完成同样的事情。
\end{itemize}

\section{示例：高效的 ASCII 字符串类型}

下面是 \texttt{Ascii} 的定义，它是一种能确保其内容始终为有效 ASCII 的字符串类型。这种类型使用了不安全特性来提供到 \texttt{String} 的零成本转换：

\begin{Verbatim}[commandchars=\\\{\},breaklines=true,breakanywhere=true,fontsize=\codefont,xleftmargin=17bp]
\PYG{k}{mod}\PYG{+w}{ }\PYG{n+nn}{my\PYGZus{}ascii}\PYG{+w}{ }\PYG{p}{\PYGZob{}}
\PYG{+w}{    }\PYG{l+s+sd}{/// 一个ASCII编码的字符串}
\PYG{+w}{    }\PYG{c+cp}{\PYGZsh{}[}\PYG{c+cp}{derive(Debug, Eq, PartialEq)}\PYG{c+cp}{]}
\PYG{+w}{    }\PYG{k}{pub}\PYG{+w}{ }\PYG{k}{struct}\PYG{+w}{ }\PYG{n+nc}{Ascii}\PYG{p}{(}
\PYG{+w}{        }\PYG{c+c1}{// 必须只持有格式良好的ASCII文本：字节范围从`0`到`0x7f`}
\PYG{+w}{        }\PYG{n+nb}{Vec}\PYG{o}{\PYGZlt{}}\PYG{k+kt}{u8}\PYG{o}{\PYGZgt{}}
\PYG{+w}{    }\PYG{p}{)}\PYG{p}{;}

\PYG{+w}{    }\PYG{k}{impl}\PYG{+w}{ }\PYG{n}{Ascii}\PYG{+w}{ }\PYG{p}{\PYGZob{}}
\PYG{+w}{        }\PYG{l+s+sd}{/// 从`bytes`的ASCII文本中创建`Ascii`。如果`bytes`包含}
\PYG{+w}{        }\PYG{l+s+sd}{/// 任何非ASCII字符，则返回`NotAsciiError`错误}
\PYG{+w}{        }\PYG{k}{pub}\PYG{+w}{ }\PYG{k}{fn}\PYG{+w}{ }\PYG{n+nf}{from\PYGZus{}bytes}\PYG{p}{(}\PYG{n}{bytes}\PYG{p}{:}\PYG{+w}{ }\PYG{n+nb}{Vec}\PYG{o}{\PYGZlt{}}\PYG{k+kt}{u8}\PYG{o}{\PYGZgt{}}\PYG{p}{)}\PYG{+w}{ }\PYG{p}{\PYGZhy{}\PYGZgt{}}\PYG{+w}{ }\PYG{n+nb}{Result}\PYG{o}{\PYGZlt{}}\PYG{n}{Ascii}\PYG{p}{,}\PYG{+w}{ }\PYG{n}{NotAsciiError}\PYG{o}{\PYGZgt{}}\PYG{+w}{ }\PYG{p}{\PYGZob{}}
\PYG{+w}{            }\PYG{k}{if}\PYG{+w}{ }\PYG{n}{bytes}\PYG{p}{.}\PYG{n}{iter}\PYG{p}{(}\PYG{p}{)}\PYG{p}{.}\PYG{n}{any}\PYG{p}{(}\PYG{o}{|}\PYG{o}{\PYGZam{}}\PYG{n}{byte}\PYG{o}{|}\PYG{+w}{ }\PYG{o}{!}\PYG{n}{byte}\PYG{p}{.}\PYG{n}{is\PYGZus{}ascii}\PYG{p}{(}\PYG{p}{)}\PYG{p}{)}\PYG{+w}{ }\PYG{p}{\PYGZob{}}
\PYG{+w}{                }\PYG{k}{return}\PYG{+w}{ }\PYG{n+nb}{Err}\PYG{p}{(}\PYG{n}{NotAsciiError}\PYG{p}{(}\PYG{n}{bytes}\PYG{p}{)}\PYG{p}{)}\PYG{p}{;}
\PYG{+w}{            }\PYG{p}{\PYGZcb{}}
\PYG{+w}{            }\PYG{n+nb}{Ok}\PYG{p}{(}\PYG{n}{Ascii}\PYG{p}{(}\PYG{n}{bytes}\PYG{p}{)}\PYG{p}{)}
\PYG{+w}{        }\PYG{p}{\PYGZcb{}}
\PYG{+w}{    }\PYG{p}{\PYGZcb{}}

\PYG{+w}{    }\PYG{c+c1}{// 当转换失败时，给出无法转换的向量。这会实现}
\PYG{+w}{    }\PYG{c+c1}{// `std::error::Error`，为保持简洁已省略}
\PYG{+w}{    }\PYG{c+cp}{\PYGZsh{}[}\PYG{c+cp}{derive(Debug, Eq, PartialEq)}\PYG{c+cp}{]}
\PYG{+w}{    }\PYG{k}{pub}\PYG{+w}{ }\PYG{k}{struct}\PYG{+w}{ }\PYG{n+nc}{NotAsciiError}\PYG{p}{(}\PYG{k}{pub}\PYG{+w}{ }\PYG{n+nb}{Vec}\PYG{o}{\PYGZlt{}}\PYG{k+kt}{u8}\PYG{o}{\PYGZgt{}}\PYG{p}{)}\PYG{p}{;}

\PYG{+w}{    }\PYG{c+c1}{// 使用不安全代码实现的安全、高效的转换}
\PYG{+w}{    }\PYG{k}{impl}\PYG{+w}{ }\PYG{n+nb}{From}\PYG{o}{\PYGZlt{}}\PYG{n}{Ascii}\PYG{o}{\PYGZgt{}}\PYG{+w}{ }\PYG{k}{for}\PYG{+w}{ }\PYG{n+nb}{String}\PYG{+w}{ }\PYG{p}{\PYGZob{}}
\PYG{+w}{        }\PYG{k}{fn}\PYG{+w}{ }\PYG{n+nf}{from}\PYG{p}{(}\PYG{n}{ascii}\PYG{p}{:}\PYG{+w}{ }\PYG{n+nc}{Ascii}\PYG{p}{)}\PYG{+w}{ }\PYG{p}{\PYGZhy{}\PYGZgt{}}\PYG{+w}{ }\PYG{n+nb}{String}\PYG{+w}{ }\PYG{p}{\PYGZob{}}
\PYG{+w}{            }\PYG{c+c1}{// 如果此模块没有bug，这就是安全的，因为格式}
\PYG{+w}{            }\PYG{c+c1}{// 良好的ASCII文本必然是格式良好的UTF\PYGZhy{}8}
\PYG{+w}{            }\PYG{k}{unsafe}\PYG{+w}{ }\PYG{p}{\PYGZob{}}\PYG{+w}{ }\PYG{n+nb}{String}\PYG{p}{::}\PYG{n}{from\PYGZus{}utf8\PYGZus{}unchecked}\PYG{p}{(}\PYG{n}{ascii}\PYG{p}{.}\PYG{l+m+mi}{0}\PYG{p}{)}\PYG{+w}{ }\PYG{p}{\PYGZcb{}}
\PYG{+w}{        }\PYG{p}{\PYGZcb{}}
\PYG{+w}{    }\PYG{p}{\PYGZcb{}}
\PYG{+w}{    }\PYG{o}{..}\PYG{p}{.}
\PYG{p}{\PYGZcb{}}
\end{Verbatim}

这个模块的关键是 \texttt{Ascii} 类型的定义。该类型本身是被标记为 \texttt{pub} 的，以令其在 \texttt{my\_ascii} 模块之外可见。但是该类型的 \texttt{Vec<u8>} 元素\textbf{不是}公共的，因此只有 \texttt{my\_ascii} 模块可以构造 \texttt{Ascii} 值或引用其元素。这使得模块的代码可以完全控制允许出现或不允许出现的内容。只要公共构造函数和方法能确保新创建的 \texttt{Ascii} 值是格式良好的并在其整个生命周期中都是如此，程序的其余部分就不会违反该规则。事实上，公共构造函数 \texttt{Ascii::from\_bytes} 在同意从给定的向量中构造 \texttt{Ascii} 之前会仔细检查它。为简洁起见，我们没有展示任何方法，但你可以想象有一组文本处理方法，并确保 \texttt{Ascii} 值始终包含正确的 ASCII 文本，就像 \texttt{String} 的方法会确保其内容始终是格式良好的 UTF-8 一样。

这种安排让我们可以非常高效地为 \texttt{String} 实现 \texttt{From<Ascii>}。不安全函数 \texttt{String::from\_utf8\_unchecked} 会获取字节向量并从中构建一个 \texttt{String}，而不会检查其内容是否为格式良好的 UTF-8 文本，该函数的契约要求其调用者对此负责。幸运的是，\texttt{Ascii} 类型强制执行的规则正是应该满足 \texttt{from\_utf8\_unchecked} 契约的规则。正如 17.2 节所解释的那样，任何 ASCII 文本块也是格式良好的 UTF-8，因此 \texttt{Ascii} 的底层 \texttt{Vec<u8>} 可以立即用作 \texttt{String} 的缓冲区。

有了这些定义，便可以这样写：

\begin{Verbatim}[commandchars=\\\{\},breaklines=true,breakanywhere=true,fontsize=\codefont,xleftmargin=17bp]
\PYG{k}{use}\PYG{+w}{ }\PYG{n}{my\PYGZus{}ascii}\PYG{p}{::}\PYG{n}{Ascii}\PYG{p}{;}

\PYG{k+kd}{let}\PYG{+w}{ }\PYG{n}{bytes}\PYG{p}{:}\PYG{+w}{ }\PYG{n+nb}{Vec}\PYG{o}{\PYGZlt{}}\PYG{k+kt}{u8}\PYG{o}{\PYGZgt{}}\PYG{+w}{ }\PYG{o}{=}\PYG{+w}{ }\PYG{l+s}{b\PYGZdq{}}\PYG{l+s}{ASCII and ye shall receive}\PYG{l+s}{\PYGZdq{}}\PYG{p}{.}\PYG{n}{to\PYGZus{}vec}\PYG{p}{(}\PYG{p}{)}\PYG{p}{;}

\PYG{c+c1}{// 这个调用不需要分配内存或复制文本，只需做扫描}
\PYG{k+kd}{let}\PYG{+w}{ }\PYG{n}{ascii}\PYG{p}{:}\PYG{+w}{ }\PYG{n+nc}{Ascii}\PYG{+w}{ }\PYG{o}{=}\PYG{+w}{ }\PYG{n}{Ascii}\PYG{p}{::}\PYG{n}{from\PYGZus{}bytes}\PYG{p}{(}\PYG{n}{bytes}\PYG{p}{)}
\PYG{+w}{    }\PYG{p}{.}\PYG{n}{unwrap}\PYG{p}{(}\PYG{p}{)}\PYG{p}{;}\PYG{+w}{ }\PYG{c+c1}{// 我们知道所选的这些字节肯定是正确的}

\PYG{c+c1}{// 这个调用是零开销的：无须分配内存、复制文本或扫描}
\PYG{k+kd}{let}\PYG{+w}{ }\PYG{n}{string}\PYG{+w}{ }\PYG{o}{=}\PYG{+w}{ }\PYG{n+nb}{String}\PYG{p}{::}\PYG{n}{from}\PYG{p}{(}\PYG{n}{ascii}\PYG{p}{)}\PYG{p}{;}

\PYG{n+nf+fm}{assert\PYGZus{}eq!}\PYG{p}{(}\PYG{n}{string}\PYG{p}{,}\PYG{+w}{ }\PYG{l+s}{\PYGZdq{}}\PYG{l+s}{ASCII and ye shall receive}\PYG{l+s}{\PYGZdq{}}\PYG{p}{)}\PYG{p}{;}
\end{Verbatim}

使用 \texttt{Ascii} 时不需要 \texttt{unsafe} 块。我们已经使用不安全操作实现了一个安全接口，并准备好仅依赖模块自己的代码而不必靠其用户的行为来满足它们的契约。

\texttt{Ascii} 只不过是 \texttt{Vec<u8>} 的包装器，但隐藏在对其内容实施额外规则的模块中。这种类型称为 newtype，这是 Rust 中的一种常见模式。Rust 自己的 \texttt{String} 类型以完全相同的方式定义，不过它的内容被限制为 UTF-8，而不是 ASCII。事实上，标准库中 \texttt{String} 的定义是这样的：

\begin{Verbatim}[commandchars=\\\{\},breaklines=true,breakanywhere=true,fontsize=\codefont,xleftmargin=17bp]
\PYG{k}{pub}\PYG{+w}{ }\PYG{k}{struct}\PYG{+w}{ }\PYG{n+nb}{String}\PYG{+w}{ }\PYG{p}{\PYGZob{}}
\PYG{+w}{    }\PYG{n}{vec}\PYG{p}{:}\PYG{+w}{ }\PYG{n+nb}{Vec}\PYG{o}{\PYGZlt{}}\PYG{k+kt}{u8}\PYG{o}{\PYGZgt{}}\PYG{p}{,}
\PYG{p}{\PYGZcb{}}
\end{Verbatim}

在机器层面，由于不认识 Rust 的类型，newtype 及其元素在内存中具有相同的表示，因此构造 newtype 根本不需要任何机器指令。在 \texttt{Ascii::from\_bytes} 中，表达式 \texttt{Ascii(bytes)} 被简单地看作 \texttt{Vec<u8>} 的一种表观，只是它现在持有一个 \texttt{Ascii} 值。同理，\texttt{String::from\_utf8\_unchecked} 在内联时可能也不需要机器指令，因为 \texttt{Vec<u8>} 现在直接作为 \texttt{String} 使用。

\section{不安全函数}

\texttt{unsafe} 函数看起来就像前面加了 \texttt{unsafe} 关键字的普通函数。\texttt{unsafe} 函数的主体自动被视为 \texttt{unsafe} 块。

只能在 \texttt{unsafe} 块中调用 \texttt{unsafe} 函数。这意味着将函数标记为 \texttt{unsafe} 会警告其调用者，为避免未定义行为，该函数具有他们必须满足的契约。

例如，下面是前面介绍的 \texttt{Ascii} 类型的新构造函数，它会从字节向量构建 \texttt{Ascii}，而不检查其内容是否为有效的 ASCII：

\begin{Verbatim}[commandchars=\\\{\},breaklines=true,breakanywhere=true,fontsize=\codefont,xleftmargin=17bp]
\PYG{c+c1}{// 以下代码必须放在`my\PYGZus{}ascii`模块内部}
\PYG{k}{impl}\PYG{+w}{ }\PYG{n}{Ascii}\PYG{+w}{ }\PYG{p}{\PYGZob{}}
\PYG{+w}{    }\PYG{l+s+sd}{/// 从`bytes`构造`Ascii`值，不检查`bytes`中是否真正包含格式良好的ASCII}
\PYG{+w}{    }\PYG{l+s+sd}{///}
\PYG{+w}{    }\PYG{l+s+sd}{/// 这个构造函数是不会出错的，它会直接返回`Ascii`，而不会像}
\PYG{+w}{    }\PYG{l+s+sd}{/// `from\PYGZus{}bytes`那样返回`Result\PYGZlt{}Ascii, NotAsciiError\PYGZgt{}`}
\PYG{+w}{    }\PYG{l+s+sd}{///}
\PYG{+w}{    }\PYG{l+s+sd}{/// \PYGZsh{} 安全性}
\PYG{+w}{    }\PYG{l+s+sd}{///}
\PYG{+w}{    }\PYG{l+s+sd}{/// 调用者必须确保`bytes`只包含ASCII字符：各字节}
\PYG{+w}{    }\PYG{l+s+sd}{/// 都不大于0x7f。否则，其行为就是未定义的}
\PYG{+w}{    }\PYG{k}{pub}\PYG{+w}{ }\PYG{k}{unsafe}\PYG{+w}{ }\PYG{k}{fn}\PYG{+w}{ }\PYG{n+nf}{from\PYGZus{}bytes\PYGZus{}unchecked}\PYG{p}{(}\PYG{n}{bytes}\PYG{p}{:}\PYG{+w}{ }\PYG{n+nb}{Vec}\PYG{o}{\PYGZlt{}}\PYG{k+kt}{u8}\PYG{o}{\PYGZgt{}}\PYG{p}{)}\PYG{+w}{ }\PYG{p}{\PYGZhy{}\PYGZgt{}}\PYG{+w}{ }\PYG{n+nc}{Ascii}\PYG{+w}{ }\PYG{p}{\PYGZob{}}
\PYG{+w}{        }\PYG{n}{Ascii}\PYG{p}{(}\PYG{n}{bytes}\PYG{p}{)}
\PYG{+w}{    }\PYG{p}{\PYGZcb{}}
\PYG{p}{\PYGZcb{}}
\end{Verbatim}

调用 \texttt{Ascii::from\_bytes\_unchecked} 的代码大概已经以某种方式知道了自己手中的向量只会包含 ASCII 字符，因此 \texttt{Ascii::from\_bytes} 坚持要执行的检查只是浪费时间，调用者也将不得不编写代码来处理他知道永远不会发生的 \texttt{Err} 结果。\texttt{Ascii::from\_bytes\_unchecked} 能让这样的调用者回避检查和错误处理。

但早些时候，为了确保 \texttt{Ascii} 值是格式良好的，我们强调了 \texttt{Ascii} 的公共构造函数和方法的重要性。\texttt{from\_bytes\_unchecked} 难道不能履行这一责任吗？

并非如此，其实 \texttt{from\_bytes\_unchecked} 通过它的契约将这些义务推脱给了调用者。这个契约的存在使得将这个函数标记为 \texttt{unsafe} 是正确的：虽然函数本身没有执行任何不安全操作，但它的调用者必须遵守某些不能靠 Rust 自动执行的规则来避免未定义行为。

真的可以通过破坏 \texttt{Ascii::from\_bytes\_unchecked} 的契约来导致未定义行为吗？是的。可以构造一个包含格式错误的 UTF-8 的 \texttt{String}，如下所示：

\begin{Verbatim}[commandchars=\\\{\},breaklines=true,breakanywhere=true,fontsize=\codefont,xleftmargin=17bp]
\PYG{c+c1}{// 将这个向量想象成用来生成ASCII的一些复杂过程的结果。但这里有问题！}
\PYG{k+kd}{let}\PYG{+w}{ }\PYG{n}{bytes}\PYG{+w}{ }\PYG{o}{=}\PYG{+w}{ }\PYG{n+nf+fm}{vec!}\PYG{p}{[}\PYG{l+m+mh}{0xf7}\PYG{p}{,}\PYG{+w}{ }\PYG{l+m+mh}{0xbf}\PYG{p}{,}\PYG{+w}{ }\PYG{l+m+mh}{0xbf}\PYG{p}{,}\PYG{+w}{ }\PYG{l+m+mh}{0xbf}\PYG{p}{]}\PYG{p}{;}

\PYG{k+kd}{let}\PYG{+w}{ }\PYG{n}{ascii}\PYG{+w}{ }\PYG{o}{=}\PYG{+w}{ }\PYG{k}{unsafe}\PYG{+w}{ }\PYG{p}{\PYGZob{}}
\PYG{+w}{    }\PYG{c+c1}{// 如果`bytes`中存有非ASCII字节，就违反了这个不安全函数的契约}
\PYG{+w}{    }\PYG{n}{Ascii}\PYG{p}{::}\PYG{n}{from\PYGZus{}bytes\PYGZus{}unchecked}\PYG{p}{(}\PYG{n}{bytes}\PYG{p}{)}
\PYG{p}{\PYGZcb{}}\PYG{p}{;}

\PYG{k+kd}{let}\PYG{+w}{ }\PYG{n}{bogus}\PYG{p}{:}\PYG{+w}{ }\PYG{n+nb}{String}\PYG{+w}{ }\PYG{o}{=}\PYG{+w}{ }\PYG{n}{ascii}\PYG{p}{.}\PYG{n}{into}\PYG{p}{(}\PYG{p}{)}\PYG{p}{;}

\PYG{c+c1}{// `bogus`现在持有格式错误的UTF\PYGZhy{}8。解析其第一个字符会生成一个不是有效Unicode}
\PYG{c+c1}{// 码点的`char`。这是未定义行为，所以语言无法说明这个断言应该是什么样的行为}
\PYG{n+nf+fm}{assert\PYGZus{}eq!}\PYG{p}{(}\PYG{n}{bogus}\PYG{p}{.}\PYG{n}{chars}\PYG{p}{(}\PYG{p}{)}\PYG{p}{.}\PYG{n}{next}\PYG{p}{(}\PYG{p}{)}\PYG{p}{.}\PYG{n}{unwrap}\PYG{p}{(}\PYG{p}{)}\PYG{+w}{ }\PYG{k}{as}\PYG{+w}{ }\PYG{k+kt}{u32}\PYG{p}{,}\PYG{+w}{ }\PYG{l+m+mh}{0x1fffff}\PYG{p}{)}\PYG{p}{;}
\end{Verbatim}

在某些平台上某些版本的 Rust 中，会观察到此断言失败并显示以下有趣的错误消息：

\begin{Verbatim}[commandchars=\\\{\},breaklines=true,breakanywhere=true,fontsize=\codefont,xleftmargin=17bp]
\PYG{n}{thread}\PYG{+w}{ }\PYG{o}{\PYGZsq{}}\PYG{n+na}{main}\PYG{o}{\PYGZsq{}}\PYG{+w}{ }\PYG{n}{panicked}\PYG{+w}{ }\PYG{n}{at}\PYG{+w}{ }\PYG{o}{\PYGZsq{}}\PYG{n+na}{assertion}\PYG{+w}{ }\PYG{n}{failed}\PYG{p}{:}\PYG{+w}{ }\PYG{err}{`}\PYG{p}{(}\PYG{n}{left}\PYG{+w}{ }\PYG{o}{=}\PYG{o}{=}\PYG{+w}{ }\PYG{n}{right}\PYG{p}{)}\PYG{err}{`}
\PYG{+w}{  }\PYG{n}{left}\PYG{p}{:}\PYG{+w}{ }\PYG{err}{`}\PYG{l+m+mi}{2097151}\PYG{err}{`}\PYG{p}{,}
\PYG{+w}{ }\PYG{n}{right}\PYG{p}{:}\PYG{+w}{ }\PYG{err}{`}\PYG{l+m+mi}{2097151}\PYG{err}{`}\PYG{o}{\PYGZsq{}}\PYG{p}{,}\PYG{+w}{ }\PYG{n}{src}\PYG{o}{/}\PYG{n}{main}\PYG{p}{.}\PYG{n}{rs}\PYG{p}{:}\PYG{l+m+mi}{42}\PYG{p}{:}\PYG{l+m+mi}{5}
\end{Verbatim}

这两个数值在我们看来明明是相等的——这不是 Rust 的错，而是前一个 \texttt{unsafe} 块所导致的。当我们说未定义行为会导致无法预测的结果时，就是这个意思。

这个例子说明了关于 bug 和不安全代码的两个关键事实。

\begin{itemize}
\item \textbf{在 \texttt{unsafe} 块之前发生的 bug 可能会破坏契约}。\texttt{unsafe} 块是否会导致未定义行为不仅取决于块本身的代码，还取决于为其提供操作目标的代码。\texttt{unsafe} 代码为满足契约所依赖的一切都与安全有关。仅当模块的其余部分都能正确维护 \texttt{Ascii} 的不变条件时，基于 \texttt{String::from\_utf8\_unchecked} 从 \texttt{Ascii} 到 \texttt{String} 的转换才是有明确定义的。
\item \textbf{离开 \texttt{unsafe} 区块后，仍可能出现此处违约的后果}。由于没有遵守不安全特性的契约而招致的未定义行为通常并不会发生在 \texttt{unsafe} 块内部。如前所述，伪造 \texttt{String} 的行为可能直到程序执行了很久之后才引发问题。
\end{itemize}

本质上，Rust 的类型检查器、借用检查器和其他静态检查都是在检查你的程序并试图构建出证据，证明它不会表现出未定义行为。如果 Rust 能成功编译程序，那么就意味着它成功地证明了你的代码是正确的。而 \texttt{unsafe} 块是这个证明中的一个缺口，也就是说，\texttt{unsafe} 块就相当于你对 Rust 说：“这段代码很好，请相信我。”你的声明正确与否可能取决于程序中会影响到此 \texttt{unsafe} 块的任意部分，并且其错误的后果也可能会出现在受此 \texttt{unsafe} 块影响的任意地点。写出 \texttt{unsafe} 关键字，就相当于你在提醒自己没能充分利用该语言的安全检查。

如果可以选择，你自然更喜欢创建不需要契约的安全接口。这些接口更容易使用，因为用户可以依靠 Rust 的安全检查来确保他们的代码没有未定义行为。即使你的实现使用了不安全特性，最好还是使用 Rust 的类型、生命周期和模块系统来满足它们的契约，同时最好只使用你能自行担保的特性，而不是把责任转嫁给你的调用者。

不过遗憾的是，在实际开发中遇到不安全函数的情况并不少见，而这些函数的文档并没有认真地解释过它们的契约。因此，你要根据自己的经验和对代码行为方式的了解自行推断出规则。如果你曾焦虑不安地想知道用 C API 或 C++ API 所做的事情是否正常，那么对这种感觉肯定也感同身受。

\section{不安全块还是不安全函数}

你可能想知道应该使用 \texttt{unsafe} 块还是将整个函数都标记为 \texttt{unsafe}。我们推荐的方法是先对函数做一些判定。

\begin{itemize}
\item 如果能正常编译，但仍可能以导致未定义行为的方式滥用函数，则必须将其标记为不安全。正确使用函数的规则是它的契约，契约的存在意味着函数是不安全的。
\item 否则，函数就是安全的。也就是说，对函数的任何类型良好的调用都不会导致未定义行为。这样的函数不应该标记为 \texttt{unsafe}。
\end{itemize}

函数是否在函数体中使用了不安全特性无关紧要，重要的是契约存在与否。之前，我们曾展示过一个没有使用不安全特性的不安全函数，以及一个使用了不安全特性的安全函数。

不要仅仅因为函数体中使用了不安全特性就把安全的函数标记为 \texttt{unsafe}。这会让函数更难使用，并使那些期望在某处找到契约说明的读者感到困惑（只要是 \texttt{unsafe} 就理当有契约说明）。相反，应该使用 \texttt{unsafe} 块，即便整个函数体只有这一个块。

\section{未定义行为}

如本章开头所述，术语\textbf{未定义行为}的意思是“Rust 坚定地认为你的代码永远不会出现的行为”。这是一个奇怪的措辞，特别是因为我们从使用其他语言的经验中就能知道这些行为\textbf{确实}会以某种频率偶然发生。为什么这个概念对厘清不安全代码的责任有帮助？

编译器是从一种编程语言到另一种编程语言的翻译器。Rust 编译器是将 Rust 程序翻译成等效的机器语言的程序。但是，说“两个使用完全不同语言的程序是等效的”意味着什么呢？

幸运的是，这个问题对程序员来说比对语言学家更容易理解。我们通常说两个程序是等效的，意思是它们在执行时总是具有相同的可见行为，比如会进行相同的系统调用，以等效的方式与外部库交互，等等。这有点儿像程序界的图灵测试：如果不能分辨自己是在与原文交互还是与译文交互，那它们就是等效的。

现在考虑以下代码：

\begin{Verbatim}[commandchars=\\\{\},breaklines=true,breakanywhere=true,fontsize=\codefont,xleftmargin=17bp]
\PYG{k+kd}{let}\PYG{+w}{ }\PYG{n}{i}\PYG{+w}{ }\PYG{o}{=}\PYG{+w}{ }\PYG{l+m+mi}{10}\PYG{p}{;}
\PYG{n}{very\PYGZus{}trustworthy}\PYG{p}{(}\PYG{o}{\PYGZam{}}\PYG{n}{i}\PYG{p}{)}\PYG{p}{;}
\PYG{n+nf+fm}{println!}\PYG{p}{(}\PYG{l+s}{\PYGZdq{}}\PYG{l+s}{\PYGZob{}\PYGZcb{}}\PYG{l+s}{\PYGZdq{}}\PYG{p}{,}\PYG{+w}{ }\PYG{n}{i}\PYG{+w}{ }\PYG{o}{*}\PYG{+w}{ }\PYG{l+m+mi}{100}\PYG{p}{)}\PYG{p}{;}
\end{Verbatim}

即使对 \texttt{very\_trustworthy} 的定义一无所知，也可以看到它仅接收对 \texttt{i} 的共享引用，因此该调用肯定无法更改 \texttt{i} 的值。由于传给 \texttt{println!} 的值永远是 \texttt{1000}，因此 Rust 可以将这段代码翻译成机器语言，就像我们写过的一样：

\begin{Verbatim}[commandchars=\\\{\},breaklines=true,breakanywhere=true,fontsize=\codefont,xleftmargin=17bp]
\PYG{n}{very\PYGZus{}trustworthy}\PYG{p}{(}\PYG{o}{\PYGZam{}}\PYG{l+m+mi}{10}\PYG{p}{)}\PYG{p}{;}
\PYG{n+nf+fm}{println!}\PYG{p}{(}\PYG{l+s}{\PYGZdq{}}\PYG{l+s}{\PYGZob{}\PYGZcb{}}\PYG{l+s}{\PYGZdq{}}\PYG{p}{,}\PYG{+w}{ }\PYG{l+m+mi}{1000}\PYG{p}{)}\PYG{p}{;}
\end{Verbatim}

这个转换后的版本与原始版本具有相同的可见行为，而且速度可能更快一点儿。但只有在保证它与原始版本具有相同含义的前提下，才值得去考虑此版本的性能。如果 \texttt{very\_trustworthy} 的定义是下面这样的该怎么办？

\begin{Verbatim}[commandchars=\\\{\},breaklines=true,breakanywhere=true,fontsize=\codefont,xleftmargin=17bp]
\PYG{k}{fn}\PYG{+w}{ }\PYG{n+nf}{very\PYGZus{}trustworthy}\PYG{p}{(}\PYG{n}{shared}\PYG{p}{:}\PYG{+w}{ }\PYG{k+kp}{\PYGZam{}}\PYG{k+kt}{i32}\PYG{p}{)}\PYG{+w}{ }\PYG{p}{\PYGZob{}}
\PYG{+w}{    }\PYG{k}{unsafe}\PYG{+w}{ }\PYG{p}{\PYGZob{}}
\PYG{+w}{        }\PYG{c+c1}{// 把共享引用转换成可变指针}
\PYG{+w}{        }\PYG{c+c1}{// 这是一种未定义行为}
\PYG{+w}{        }\PYG{k+kd}{let}\PYG{+w}{ }\PYG{n}{mutable}\PYG{+w}{ }\PYG{o}{=}\PYG{+w}{ }\PYG{n}{shared}\PYG{+w}{ }\PYG{k}{as}\PYG{+w}{ }\PYG{o}{*}\PYG{k}{const}\PYG{+w}{ }\PYG{k+kt}{i32}\PYG{+w}{ }\PYG{k}{as}\PYG{+w}{ }\PYG{o}{*}\PYG{k}{mut}\PYG{+w}{ }\PYG{k+kt}{i32}\PYG{p}{;}
\PYG{+w}{        }\PYG{o}{*}\PYG{n}{mutable}\PYG{+w}{ }\PYG{o}{=}\PYG{+w}{ }\PYG{l+m+mi}{20}\PYG{p}{;}
\PYG{+w}{    }\PYG{p}{\PYGZcb{}}
\PYG{p}{\PYGZcb{}}
\end{Verbatim}

这段代码打破了共享引用的规则：它将 \texttt{i} 的值更改为 \texttt{20}，即使这个值应该被冻结（因为 \texttt{i} 是为了共享而借用的）。结果，我们对调用者所做的转换现在有了非常明显的效果：如果 \texttt{Rust} 转换这段代码，那么程序就会打印 \texttt{1000}；如果它保留代码并使用 \texttt{i} 的新值，则程序会打印 \texttt{2000}。在 \texttt{very\_trustworthy} 中打破共享引用的规则意味着共享引用在其调用者中的行为可能不再符合预期了。

这种问题几乎会出现在 Rust 可能尝试的每一种转换中，其中就包括：即使把一个函数内联到它的调用点，也仍然可以假设当被调用者完成时，控制流就会返回到调用点。然而我们在本章开头就给出过一个违反了该假设的问题代码示例。

Rust（或任何其他语言）基本上不可能评估出对程序的转换是否保留了其含义，除非它可以相信语言的基本特性会按原本的设计运行。一段代码是否可信，不仅取决于手头的这部分代码，还取决于程序中隔得比较远的其他部分的代码。为了对代码做任何处理，Rust 必须假设你的程序的其余部分是“遵纪守法”的。

下面是 Rust 判断程序是否“遵纪守法”的规则。

\begin{itemize}
\item 程序不得读取未初始化的内存。
\item 程序不得创建无效的原始值。

\begin{itemize}
\item 引用、Box 值或 \texttt{fn} 指针为空（\texttt{null}）。
\item \texttt{bool} 值非 \texttt{0} 且非 \texttt{1}。
\item \texttt{enum} 值具有无效判别值。
\item \texttt{char} 值无效，比如存在半代用区的 Unicode 码点。
\item \texttt{str} 值不是格式良好的 UTF-8。
\item 胖指针具有无效虚表或 slice 长度。
\item \texttt{never} 类型的任何值，可以写作 \texttt{!}，只能用于不会返回的函数。
\end{itemize}

\item 程序必须遵守第 5 章中解释过的引用规则。任何引用的生命周期都不能超出其引用目标，共享访问是只读访问，可变访问是独占访问。
\item 程序不得对空指针、未正确对齐的指针或悬空指针进行解引用。
\item 程序不得使用指针访问与此指针关联的分配区之外的内存。22.8.1 节会详细解释此规则。
\item 程序必须没有数据竞争。当两个线程在没有同步保护的情况下访问同一个内存位置，并且至少有一个访问是写入时，就会发生数据竞争。
\item 程序不得对借助外部函数接口进行的跨语言调用进行栈展开（参见 7.1.1 节）。
\item 程序必须遵守标准库函数的契约。
\end{itemize}

由于还没有针对 \texttt{unsafe} 代码的 Rust 语义的完整模型，因此该列表可能会随着时间的推移而演变，但上述这些规则仍然有效。

任何违反这些规则的行为都构成了未定义行为，并让 Rust 试图优化你的程序并将其翻译成机器语言的努力变得不可信。如果你违反最后一条规则并将格式错误的 UTF-8 传给 \texttt{String::from\_utf8\_unchecked}，那么没准儿 2097151 真会不等于 2097151。

未使用不安全特性的 Rust 代码只要编译通过就可以保证会遵守前面的所有规则。（假设编译器自身没有 bug——我们正为之努力，但没有 bug 永远只能是个理想。）只有在使用不安全特性时，遵守这些规则才会成为你的责任。

而在 C 和 C++ 中，你的程序“在编译期没有错误或警告”这件事意义不大。正如本书前面所提到的，即使那些一直坚持着高标准且备受推崇的项目所编写的最好的 C 程序和 C++ 程序，也会在实践中表现出未定义行为。

\section{不安全特型}

\textbf{不安全特型}是一种特型，用于表示这里存在某种 Rust 无法检查也无法强制保障的契约。实现者必须满足它，以规避未定义行为。要实现不安全特型，就必须将实现标记为不安全的。你需要了解此特型的契约并确保自己的类型能满足它。

类型变量以某个不安全特型为限界的函数通常是自身使用了不安全特型的函数，并且只能依靠此不安全特型的契约来满足那些不安全特型的契约。对此特型的不正确实现可能导致这样的函数出现未定义行为。

\texttt{std::marker::Send} 和 \texttt{std::marker::Sync} 是不安全特型的典型示例。这些特型没有定义任何方法，因此你可以用喜欢的任意类型来轻松实现它们。但它们确实有契约：\texttt{Send} 要求实现者能安全地转移给另一个线程，而 \texttt{Sync} 要求实现者能安全地通过共享引用在线程之间共享。如果为不合适的类型实现了 \texttt{Send}，就会使 \texttt{std::sync::Mutex} 在数据竞争中不再安全。

举个简单的例子，Rust 标准库曾包含一个不安全特型 \texttt{core::nonzero::Zeroable}，该特型用于标记出可通过将所有字节设置为 0 来进行安全初始化的类型。显然，将 \texttt{usize} 变量归零肯定没问题，但将 \texttt{\&T} 归零就会带来一个空引用，如果解引用，则会导致崩溃。对于 \texttt{Zeroable} 的类型，可以进行一些优化：可以使用 \texttt{std::ptr::write\_bytes}（\texttt{memset} 在 Rust 中的等价物）或者用能分配全零内存页的系统调用来快速初始化数组。（\texttt{Zeroable} 是不稳定的，在 Rust 1.26 的 \texttt{num} crate 中被转移到仅供内部使用，但它是一个优秀、简单且真实的例子。）

\texttt{Zeroable} 是一个典型的标记特型，缺少方法或关联类型：

\begin{Verbatim}[commandchars=\\\{\},breaklines=true,breakanywhere=true,fontsize=\codefont,xleftmargin=17bp]
\PYG{k}{pub}\PYG{+w}{ }\PYG{k}{unsafe}\PYG{+w}{ }\PYG{k}{trait}\PYG{+w}{ }\PYG{n}{Zeroable}\PYG{+w}{ }\PYG{p}{\PYGZob{}}\PYG{p}{\PYGZcb{}}
\end{Verbatim}

它对适用类型的实现同样简单明了：

\begin{Verbatim}[commandchars=\\\{\},breaklines=true,breakanywhere=true,fontsize=\codefont,xleftmargin=17bp]
\PYG{k}{unsafe}\PYG{+w}{ }\PYG{k}{impl}\PYG{+w}{ }\PYG{n}{Zeroable}\PYG{+w}{ }\PYG{k}{for}\PYG{+w}{ }\PYG{k+kt}{u8}\PYG{+w}{ }\PYG{p}{\PYGZob{}}\PYG{p}{\PYGZcb{}}
\PYG{k}{unsafe}\PYG{+w}{ }\PYG{k}{impl}\PYG{+w}{ }\PYG{n}{Zeroable}\PYG{+w}{ }\PYG{k}{for}\PYG{+w}{ }\PYG{k+kt}{i32}\PYG{+w}{ }\PYG{p}{\PYGZob{}}\PYG{p}{\PYGZcb{}}
\PYG{k}{unsafe}\PYG{+w}{ }\PYG{k}{impl}\PYG{+w}{ }\PYG{n}{Zeroable}\PYG{+w}{ }\PYG{k}{for}\PYG{+w}{ }\PYG{k+kt}{usize}\PYG{+w}{ }\PYG{p}{\PYGZob{}}\PYG{p}{\PYGZcb{}}
\PYG{c+c1}{// 以及所有整数类型}
\end{Verbatim}

有了这些定义，就可以编写一个函数来快速分配给定长度的包含 \texttt{Zeroable} 类型的向量了：

\begin{Verbatim}[commandchars=\\\{\},breaklines=true,breakanywhere=true,fontsize=\codefont,xleftmargin=17bp]
\PYG{k}{use}\PYG{+w}{ }\PYG{n}{core}\PYG{p}{::}\PYG{n}{nonzero}\PYG{p}{::}\PYG{n}{Zeroable}\PYG{p}{;}

\PYG{k}{fn}\PYG{+w}{ }\PYG{n+nf}{zeroed\PYGZus{}vector}\PYG{o}{\PYGZlt{}}\PYG{n}{T}\PYG{o}{\PYGZgt{}}\PYG{p}{(}\PYG{n}{len}\PYG{p}{:}\PYG{+w}{ }\PYG{k+kt}{usize}\PYG{p}{)}\PYG{+w}{ }\PYG{p}{\PYGZhy{}\PYGZgt{}}\PYG{+w}{ }\PYG{n+nb}{Vec}\PYG{o}{\PYGZlt{}}\PYG{n}{T}\PYG{o}{\PYGZgt{}}
\PYG{+w}{    }\PYG{k}{where}\PYG{+w}{ }\PYG{n}{T}\PYG{p}{:}\PYG{+w}{ }\PYG{n+nc}{Zeroable}
\PYG{p}{\PYGZob{}}
\PYG{+w}{    }\PYG{k+kd}{let}\PYG{+w}{ }\PYG{k}{mut}\PYG{+w}{ }\PYG{n}{vec}\PYG{+w}{ }\PYG{o}{=}\PYG{+w}{ }\PYG{n+nb}{Vec}\PYG{p}{::}\PYG{n}{with\PYGZus{}capacity}\PYG{p}{(}\PYG{n}{len}\PYG{p}{)}\PYG{p}{;}
\PYG{+w}{    }\PYG{k}{unsafe}\PYG{+w}{ }\PYG{p}{\PYGZob{}}
\PYG{+w}{        }\PYG{n}{std}\PYG{p}{::}\PYG{n}{ptr}\PYG{p}{::}\PYG{n}{write\PYGZus{}bytes}\PYG{p}{(}\PYG{n}{vec}\PYG{p}{.}\PYG{n}{as\PYGZus{}mut\PYGZus{}ptr}\PYG{p}{(}\PYG{p}{)}\PYG{p}{,}\PYG{+w}{ }\PYG{l+m+mi}{0}\PYG{p}{,}\PYG{+w}{ }\PYG{n}{len}\PYG{p}{)}\PYG{p}{;}
\PYG{+w}{        }\PYG{n}{vec}\PYG{p}{.}\PYG{n}{set\PYGZus{}len}\PYG{p}{(}\PYG{n}{len}\PYG{p}{)}\PYG{p}{;}
\PYG{+w}{    }\PYG{p}{\PYGZcb{}}
\PYG{+w}{    }\PYG{n}{vec}
\PYG{p}{\PYGZcb{}}
\end{Verbatim}

这个函数会首先创建一个具有所需容量的空 \texttt{Vec}，然后调用 \texttt{write\_bytes} 以用 0 填充未占用的缓冲区。（\texttt{write\_byte} 函数会将 \texttt{len} 视为 \texttt{T} 元素的数量，而不是字节数，因此该调用确实会填充整个缓冲区。）向量的 \texttt{set\_len} 方法只会更改其长度而不会对缓冲区做任何事，这是不安全的，因为必须保证新的缓冲区空间确实包含已正确初始化的 \texttt{T} 类型值。不过这正是 \texttt{T: Zeroable} 类型限界所保证的：全零的字节块表示有效的 \texttt{T} 值。我们对 \texttt{set\_len} 的使用是安全的。

下面我们来使用这个函数：

\begin{Verbatim}[commandchars=\\\{\},breaklines=true,breakanywhere=true,fontsize=\codefont,xleftmargin=17bp]
\PYG{k+kd}{let}\PYG{+w}{ }\PYG{n}{v}\PYG{p}{:}\PYG{+w}{ }\PYG{n+nb}{Vec}\PYG{o}{\PYGZlt{}}\PYG{k+kt}{usize}\PYG{o}{\PYGZgt{}}\PYG{+w}{ }\PYG{o}{=}\PYG{+w}{ }\PYG{n}{zeroed\PYGZus{}vector}\PYG{p}{(}\PYG{l+m+mi}{100\PYGZus{}000}\PYG{p}{)}\PYG{p}{;}
\PYG{n+nf+fm}{assert!}\PYG{p}{(}\PYG{n}{v}\PYG{p}{.}\PYG{n}{iter}\PYG{p}{(}\PYG{p}{)}\PYG{p}{.}\PYG{n}{all}\PYG{p}{(}\PYG{o}{|}\PYG{o}{\PYGZam{}}\PYG{n}{u}\PYG{o}{|}\PYG{+w}{ }\PYG{n}{u}\PYG{+w}{ }\PYG{o}{=}\PYG{o}{=}\PYG{+w}{ }\PYG{l+m+mi}{0}\PYG{p}{)}\PYG{p}{)}\PYG{p}{;}
\end{Verbatim}

显然，\texttt{Zeroable} 一定是一个不安全特型，因为不遵守其契约的实现可能会导致未定义行为：

\begin{Verbatim}[commandchars=\\\{\},breaklines=true,breakanywhere=true,fontsize=\codefont,xleftmargin=17bp]
\PYG{k}{struct}\PYG{+w}{ }\PYG{n+nc}{HoldsRef}\PYG{o}{\PYGZlt{}}\PYG{o}{\PYGZsq{}}\PYG{n+na}{a}\PYG{o}{\PYGZgt{}}\PYG{p}{(}\PYG{o}{\PYGZam{}}\PYG{o}{\PYGZsq{}}\PYG{n+na}{a}\PYG{+w}{ }\PYG{k}{mut}\PYG{+w}{ }\PYG{k+kt}{i32}\PYG{p}{)}\PYG{p}{;}

\PYG{k}{unsafe}\PYG{+w}{ }\PYG{k}{impl}\PYG{o}{\PYGZlt{}}\PYG{o}{\PYGZsq{}}\PYG{n+na}{a}\PYG{o}{\PYGZgt{}}\PYG{+w}{ }\PYG{n}{Zeroable}\PYG{+w}{ }\PYG{k}{for}\PYG{+w}{ }\PYG{n}{HoldsRef}\PYG{o}{\PYGZlt{}}\PYG{o}{\PYGZsq{}}\PYG{n+na}{a}\PYG{o}{\PYGZgt{}}\PYG{+w}{ }\PYG{p}{\PYGZob{}}\PYG{+w}{ }\PYG{p}{\PYGZcb{}}

\PYG{k+kd}{let}\PYG{+w}{ }\PYG{k}{mut}\PYG{+w}{ }\PYG{n}{v}\PYG{p}{:}\PYG{+w}{ }\PYG{n+nb}{Vec}\PYG{o}{\PYGZlt{}}\PYG{n}{HoldsRef}\PYG{o}{\PYGZgt{}}\PYG{+w}{ }\PYG{o}{=}\PYG{+w}{ }\PYG{n}{zeroed\PYGZus{}vector}\PYG{p}{(}\PYG{l+m+mi}{1}\PYG{p}{)}\PYG{p}{;}
\PYG{o}{*}\PYG{n}{v}\PYG{p}{[}\PYG{l+m+mi}{0}\PYG{p}{]}\PYG{p}{.}\PYG{l+m+mi}{0}\PYG{+w}{ }\PYG{o}{=}\PYG{+w}{ }\PYG{l+m+mi}{1}\PYG{p}{;}\PYG{+w}{   }\PYG{c+c1}{// 崩溃：对空指针解引用}
\end{Verbatim}

Rust 不知道 \texttt{Zeroable} 意味着什么，所以无从判断它何时会被实现为不合适的类型。与其他任何不安全特性一样，如何理解并遵守不安全特型的契约由你来决定。

请注意，不安全代码不得依赖于普通的、安全的特型在实现上的正确性。假设有一个 \texttt{std::hash::Hasher} 特型的实现，它只会返回一个随机哈希值，与被哈希的值无关。该特型要求对一些相同的位进行两次哈希后必须生成相同的哈希值，但此实现无法满足该要求，这根本就不正确。但因为 \texttt{Hasher} 并不是不安全特型，所以不安全代码在使用这个哈希器时不得表现出未定义行为。\footnote{\texttt{Hasher} 的行为是有明确定义的，别人实现错了并不等于此特型未定义。——译者注}为了满足“可以使用不安全特性”这条契约，\texttt{std::collections::HashMap} 类型是经过精心编写的，但并未考虑哈希器自身行为出错的可能性。当然，这样一来该哈希表将无法正常运行：查找将失败，条目将随机出现和消失。但该哈希表并不存在未定义行为。

\section{裸指针}

裸指针在 Rust 中就是一种不受约束的指针。你可以使用裸指针来创建 Rust 的受检查指针类型不能创建的各种结构，比如双向链表或任意对象图。但是因为裸指针非常灵活，Rust 无法判断你是否在安全地使用它们，所以只能在 \texttt{unsafe} 块中对它们解引用。

裸指针本质上等效于 C 指针或 C++ 指针，因此在与这些语言编写的代码进行交互时它们也很有用。

裸指针有以下两种类型。

\begin{itemize}
\item \texttt{*mut T} 是指向 \texttt{T} 的允许修改其引用目标的裸指针。
\item \texttt{*const T} 是指向 \texttt{T} 的只允许读取其引用目标的裸指针。
\end{itemize}

（没有单纯的 \texttt{*T} 类型，必须始终指定 \texttt{const} 或 \texttt{mut}。）

可以把引用转换成裸指针，并使用 \texttt{*} 运算符对其解引用：

\begin{Verbatim}[commandchars=\\\{\},breaklines=true,breakanywhere=true,fontsize=\codefont,xleftmargin=17bp]
\PYG{k+kd}{let}\PYG{+w}{ }\PYG{k}{mut}\PYG{+w}{ }\PYG{n}{x}\PYG{+w}{ }\PYG{o}{=}\PYG{+w}{ }\PYG{l+m+mi}{10}\PYG{p}{;}
\PYG{k+kd}{let}\PYG{+w}{ }\PYG{n}{ptr\PYGZus{}x}\PYG{+w}{ }\PYG{o}{=}\PYG{+w}{ }\PYG{o}{\PYGZam{}}\PYG{k}{mut}\PYG{+w}{ }\PYG{n}{x}\PYG{+w}{ }\PYG{k}{as}\PYG{+w}{ }\PYG{o}{*}\PYG{k}{mut}\PYG{+w}{ }\PYG{k+kt}{i32}\PYG{p}{;}

\PYG{k+kd}{let}\PYG{+w}{ }\PYG{n}{y}\PYG{+w}{ }\PYG{o}{=}\PYG{+w}{ }\PYG{n+nb}{Box}\PYG{p}{::}\PYG{n}{new}\PYG{p}{(}\PYG{l+m+mi}{20}\PYG{p}{)}\PYG{p}{;}
\PYG{k+kd}{let}\PYG{+w}{ }\PYG{n}{ptr\PYGZus{}y}\PYG{+w}{ }\PYG{o}{=}\PYG{+w}{ }\PYG{o}{\PYGZam{}}\PYG{o}{*}\PYG{n}{y}\PYG{+w}{ }\PYG{k}{as}\PYG{+w}{ }\PYG{o}{*}\PYG{k}{const}\PYG{+w}{ }\PYG{k+kt}{i32}\PYG{p}{;}

\PYG{k}{unsafe}\PYG{+w}{ }\PYG{p}{\PYGZob{}}
\PYG{+w}{    }\PYG{o}{*}\PYG{n}{ptr\PYGZus{}x}\PYG{+w}{ }\PYG{o}{+}\PYG{o}{=}\PYG{+w}{ }\PYG{o}{*}\PYG{n}{ptr\PYGZus{}y}\PYG{p}{;}
\PYG{p}{\PYGZcb{}}
\PYG{n+nf+fm}{assert\PYGZus{}eq!}\PYG{p}{(}\PYG{n}{x}\PYG{p}{,}\PYG{+w}{ }\PYG{l+m+mi}{30}\PYG{p}{)}\PYG{p}{;}
\end{Verbatim}

与 Box 和引用不同，裸指针可以为空，就像 C 中的 \texttt{NULL} 或 C++ 中的 \texttt{nullptr}：

\begin{Verbatim}[commandchars=\\\{\},breaklines=true,breakanywhere=true,fontsize=\codefont,xleftmargin=17bp]
\PYG{k}{fn}\PYG{+w}{ }\PYG{n+nf}{option\PYGZus{}to\PYGZus{}raw}\PYG{o}{\PYGZlt{}}\PYG{n}{T}\PYG{o}{\PYGZgt{}}\PYG{p}{(}\PYG{n}{opt}\PYG{p}{:}\PYG{+w}{ }\PYG{n+nb}{Option}\PYG{o}{\PYGZlt{}}\PYG{o}{\PYGZam{}}\PYG{n}{T}\PYG{o}{\PYGZgt{}}\PYG{p}{)}\PYG{+w}{ }\PYG{p}{\PYGZhy{}\PYGZgt{}}\PYG{+w}{ }\PYG{o}{*}\PYG{k}{const}\PYG{+w}{ }\PYG{n}{T}\PYG{+w}{ }\PYG{p}{\PYGZob{}}
\PYG{+w}{    }\PYG{k}{match}\PYG{+w}{ }\PYG{n}{opt}\PYG{+w}{ }\PYG{p}{\PYGZob{}}
\PYG{+w}{        }\PYG{n+nb}{None}\PYG{+w}{ }\PYG{o}{=}\PYG{o}{\PYGZgt{}}\PYG{+w}{ }\PYG{n}{std}\PYG{p}{::}\PYG{n}{ptr}\PYG{p}{::}\PYG{n}{null}\PYG{p}{(}\PYG{p}{)}\PYG{p}{,}
\PYG{+w}{        }\PYG{n+nb}{Some}\PYG{p}{(}\PYG{n}{r}\PYG{p}{)}\PYG{+w}{ }\PYG{o}{=}\PYG{o}{\PYGZgt{}}\PYG{+w}{ }\PYG{n}{r}\PYG{+w}{ }\PYG{k}{as}\PYG{+w}{ }\PYG{o}{*}\PYG{k}{const}\PYG{+w}{ }\PYG{n}{T}
\PYG{+w}{    }\PYG{p}{\PYGZcb{}}
\PYG{p}{\PYGZcb{}}

\PYG{n+nf+fm}{assert!}\PYG{p}{(}\PYG{o}{!}\PYG{n}{option\PYGZus{}to\PYGZus{}raw}\PYG{p}{(}\PYG{n+nb}{Some}\PYG{p}{(}\PYG{o}{\PYGZam{}}\PYG{p}{(}\PYG{l+s}{\PYGZdq{}}\PYG{l+s}{pea}\PYG{l+s}{\PYGZdq{}}\PYG{p}{,}\PYG{+w}{ }\PYG{l+s}{\PYGZdq{}}\PYG{l+s}{pod}\PYG{l+s}{\PYGZdq{}}\PYG{p}{)}\PYG{p}{)}\PYG{p}{)}\PYG{p}{.}\PYG{n}{is\PYGZus{}null}\PYG{p}{(}\PYG{p}{)}\PYG{p}{)}\PYG{p}{;}
\PYG{n+nf+fm}{assert\PYGZus{}eq!}\PYG{p}{(}\PYG{n}{option\PYGZus{}to\PYGZus{}raw}\PYG{p}{:}\PYG{p}{:}\PYG{o}{\PYGZlt{}}\PYG{k+kt}{i32}\PYG{o}{\PYGZgt{}}\PYG{p}{(}\PYG{n+nb}{None}\PYG{p}{)}\PYG{p}{,}\PYG{+w}{ }\PYG{n}{std}\PYG{p}{::}\PYG{n}{ptr}\PYG{p}{::}\PYG{n}{null}\PYG{p}{(}\PYG{p}{)}\PYG{p}{)}\PYG{p}{;}
\end{Verbatim}

这个例子中没有 \texttt{unsafe} 块：创建裸指针、传递裸指针和比较裸指针都是安全的。只有解引用裸指针是不安全的。

指向无固定大小类型的裸指针是胖指针，就像相应的引用或 \texttt{Box} 类型一样。\texttt{*const [u8]} 指针包括一个长度和地址，而像 \texttt{*mut dyn std::io::Write} 指针这样的特型对象则会携带一个虚表。

尽管 Rust 会在各种情况下隐式解引用安全指针类型，但对裸指针解引用必须是显式的。

\begin{itemize}
\item \texttt{.} 运算符不会隐式解引用裸指针，必须写成 \texttt{(*raw).field} 或 \texttt{(*raw).method(...)}。
\item 裸指针没有实现 \texttt{Deref}，因此隐式解引用不适合它们。
\item \texttt{==}、\texttt{<} 等运算符将裸指针作为地址进行比较：如果两个裸指针指向内存中的相同位置，那它们就相等。类似地，对裸指针进行哈希处理会针对其指向的地址值本身，而不会针对其引用目标的值。
\item 像 \texttt{std::fmt::Display} 这样的格式化特型会自动追踪引用，但根本不会处理裸指针。\texttt{std::fmt::Debug} 和 \texttt{std::fmt::Pointer} 是例外，它们会将裸指针展示为十六进制地址，而不会解引用它们。
\end{itemize}

与 C 和 C++ 中的 \texttt{+} 运算符不同，Rust 的 \texttt{+} 不会处理裸指针，但可以通过它们的 \texttt{offset} 方法和 \texttt{wrapping\_offset} 方法或更方便的 \texttt{add} 方法、\texttt{sub} 方法、\texttt{wrapping\_add} 方法和 \texttt{wrapping\_sub} 方法执行指针运算。反过来，\texttt{offset\_from} 方法会以字节为单位求出两个指针之间的距离，不过需要确保开始和结束位于同一个内存区域，比如在同一个 \texttt{Vec} 中：

\begin{Verbatim}[commandchars=\\\{\},breaklines=true,breakanywhere=true,fontsize=\codefont,xleftmargin=17bp]
\PYG{k+kd}{let}\PYG{+w}{ }\PYG{n}{trucks}\PYG{+w}{ }\PYG{o}{=}\PYG{+w}{ }\PYG{n+nf+fm}{vec!}\PYG{p}{[}\PYG{l+s}{\PYGZdq{}}\PYG{l+s}{garbage truck}\PYG{l+s}{\PYGZdq{}}\PYG{p}{,}\PYG{+w}{ }\PYG{l+s}{\PYGZdq{}}\PYG{l+s}{dump truck}\PYG{l+s}{\PYGZdq{}}\PYG{p}{,}\PYG{+w}{ }\PYG{l+s}{\PYGZdq{}}\PYG{l+s}{moonstruck}\PYG{l+s}{\PYGZdq{}}\PYG{p}{]}\PYG{p}{;}
\PYG{k+kd}{let}\PYG{+w}{ }\PYG{n}{first}\PYG{p}{:}\PYG{+w}{ }\PYG{o}{*}\PYG{k}{const}\PYG{+w}{ }\PYG{o}{\PYGZam{}}\PYG{k+kt}{str}\PYG{+w}{ }\PYG{o}{=}\PYG{+w}{ }\PYG{o}{\PYGZam{}}\PYG{n}{trucks}\PYG{p}{[}\PYG{l+m+mi}{0}\PYG{p}{]}\PYG{p}{;}
\PYG{k+kd}{let}\PYG{+w}{ }\PYG{n}{last}\PYG{p}{:}\PYG{+w}{ }\PYG{o}{*}\PYG{k}{const}\PYG{+w}{ }\PYG{o}{\PYGZam{}}\PYG{k+kt}{str}\PYG{+w}{ }\PYG{o}{=}\PYG{+w}{ }\PYG{o}{\PYGZam{}}\PYG{n}{trucks}\PYG{p}{[}\PYG{l+m+mi}{2}\PYG{p}{]}\PYG{p}{;}
\PYG{n+nf+fm}{assert\PYGZus{}eq!}\PYG{p}{(}\PYG{k}{unsafe}\PYG{+w}{ }\PYG{p}{\PYGZob{}}\PYG{+w}{ }\PYG{n}{last}\PYG{p}{.}\PYG{n}{offset\PYGZus{}from}\PYG{p}{(}\PYG{n}{first}\PYG{p}{)}\PYG{+w}{ }\PYG{p}{\PYGZcb{}}\PYG{p}{,}\PYG{+w}{ }\PYG{l+m+mi}{2}\PYG{p}{)}\PYG{p}{;}
\PYG{n+nf+fm}{assert\PYGZus{}eq!}\PYG{p}{(}\PYG{k}{unsafe}\PYG{+w}{ }\PYG{p}{\PYGZob{}}\PYG{+w}{ }\PYG{n}{first}\PYG{p}{.}\PYG{n}{offset\PYGZus{}from}\PYG{p}{(}\PYG{n}{last}\PYG{p}{)}\PYG{+w}{ }\PYG{p}{\PYGZcb{}}\PYG{p}{,}\PYG{+w}{ }\PYG{o}{\PYGZhy{}}\PYG{l+m+mi}{2}\PYG{p}{)}\PYG{p}{;}
\end{Verbatim}

\texttt{first} 和 \texttt{last} 不需要显式转换，只需指定类型即可。Rust 会将引用隐式转换成裸指针（当然，反过来肯定不成立）。

\texttt{as} 运算符允许从引用到裸指针或两个裸指针类型之间几乎所有的合理转换。但是，可能需要将复杂的转换分解为一系列更简单的步骤。例如：

\begin{Verbatim}[commandchars=\\\{\},breaklines=true,breakanywhere=true,fontsize=\codefont,xleftmargin=17bp]
\PYG{o}{\PYGZam{}}\PYG{n+nf+fm}{vec!}\PYG{p}{[}\PYG{l+m+mi}{42\PYGZus{}}\PYG{k}{u8}\PYG{p}{]}\PYG{+w}{ }\PYG{k}{as}\PYG{+w}{ }\PYG{o}{*}\PYG{k}{const}\PYG{+w}{ }\PYG{n+nb}{String}\PYG{p}{;}\PYG{+w}{  }\PYG{c+c1}{// 错误：无效的转换}
\PYG{o}{\PYGZam{}}\PYG{n+nf+fm}{vec!}\PYG{p}{[}\PYG{l+m+mi}{42\PYGZus{}}\PYG{k}{u8}\PYG{p}{]}\PYG{+w}{ }\PYG{k}{as}\PYG{+w}{ }\PYG{o}{*}\PYG{k}{const}\PYG{+w}{ }\PYG{n+nb}{Vec}\PYG{o}{\PYGZlt{}}\PYG{k+kt}{u8}\PYG{o}{\PYGZgt{}}\PYG{+w}{ }\PYG{k}{as}\PYG{+w}{ }\PYG{o}{*}\PYG{k}{const}\PYG{+w}{ }\PYG{n+nb}{String}\PYG{p}{;}\PYG{+w}{  }\PYG{c+c1}{// 这样可以转换}
\end{Verbatim}

请注意，\texttt{as} 不会将裸指针转换为引用。这样的转换不安全，\texttt{as} 应该保持安全操作。因此，必须在 \texttt{unsafe} 块中对裸指针解引用，然后再借用其结果值。

这样操作时要非常小心：以这种方式生成的引用具有不受约束的生命周期，它可以存续多长时间没有限制，因为裸指针没有给 Rust 提供任何能做出这种决定的依据。23.5 节会展示几个如何正确限制生命周期的示例。

许多类型有 \texttt{as\_ptr} 方法和 \texttt{as\_mut\_ptr} 方法，它们会返回指向其内容的裸指针。例如，数组切片和字符串会返回指向它们第一个元素的指针，而一些迭代器会返回指向它们将生成的下一个元素的指针。像 \texttt{Box}、\texttt{Rc} 和 \texttt{Arc} 这样的拥有型指针类型都有 \texttt{into\_raw} 函数和 \texttt{from\_raw} 函数，可以与裸指针相互转换，其中一些方法的契约强加了出人意料的要求，因此在使用之前务必检查一下它们的文档。

还可以通过转换整数来构造裸指针，不过你唯一可以信任的整数通常就是从指针转换来的。22.8.2 节就以这种方式使用了裸指针。

与引用不同，裸指针既不是 \texttt{Send} 的也不是 \texttt{Sync} 的。因此，在默认情况下，任何包含裸指针的类型都不会实现这些特型。在线程之间发送或共享裸指针本身其实并没有什么不安全的，毕竟，无论它们“走”到哪里，你都需要一个 \texttt{unsafe} 块来解引用它们。但是考虑到裸指针经常扮演的角色，语言设计者认为还是现在这种默认使用方式更好。22.7 节讨论过如何自己实现 \texttt{Send} 和 \texttt{Sync}。

\subsection{安全地解引用裸指针}

以下是安全使用裸指针的一些常识性指南。

\begin{itemize}
\item 解引用空指针或悬空指针是未定义行为，引用未初始化的内存或超出作用域的值也一样。
\item 解引用未针对其引用目标的类型正确对齐的指针是未定义行为。
\item 只有在遵守了第 5 章解释过的引用安全规则（任何引用的生命周期都不能超出其引用目标，共享访问是只读访问，可变访问是独占访问）的前提下，才能从解引用的裸指针中借用值。（很容易意外违反这条规则，因为裸指针通常用于创建具有非标准共享或所有权的数据结构。）
\item 仅当引用目标是所属类型的格式良好的值时，才能使用裸指针的引用目标。例如，必须确保解引用 \texttt{*const char} 后会产生一个正确的、不在半代用区的 Unicode 码点。
\item 如果想在特定的裸指针上使用 \texttt{offset} 方法和 \texttt{wrapping\_offset} 方法，那么该裸指针只能指向原初（original）指针所引用的变量内部的字节或分配在堆上的内存块内部的字节，或者指向上述两个区域之外的第一字节。

如果通过将指针转换为整数，对整数进行运算，然后将其转换回指针的方式进行指针运算，则结果必须是 \texttt{offset} 方法的规则允许生成的指针。
\item 如果要给裸指针的引用目标赋值，则不得违反引用目标所属的任何类型的不变条件。如果你有一个 \texttt{*mut u8} 指向 \texttt{String} 中的一字节，那么在该 \texttt{u8} 中存储的值必须能让 \texttt{String} 保持为格式良好的 UTF-8。
\end{itemize}

抛开借用规则不谈，上述规则与在 C 或 C++ 中使用指针时必须遵守的规则基本上是一样的。

不得违反类型不变条件的原因应该很清楚。许多 Rust 标准库类型在其实现中使用了不安全代码，但仍然提供了安全接口，前提是 Rust 的安全检查、模块系统和可见性规则能得到遵守。使用裸指针来规避这些保护措施可能会导致未定义行为。

裸指针的完整、准确的契约不容易表述，并且可能随着语言的发展而改变。但本节概要表述的这些原则应该让你处于安全地带。

\subsection{示例：\texttt{RefWithFlag}}

下面这个例子说明了如何采用裸指针实现经典\footnote{好吧，应该说是成就我们（本书作者）的经典。}的位级 hack，并将其包装为完全安全的 Rust 类型。这个模块定义了一个类型 \texttt{RefWithFlag<'a, T>}，它同时包含一个 \texttt{\&'a T} 和一个 \texttt{bool}，就像元组 \texttt{(\&'a T, bool)} 一样，但仍然设法只占用了一个机器字而不是两个。这种技术在垃圾回收器和虚拟机中经常使用，其中某些类型（比如表示对象的类型）的数量多到就算只向每个值添加一个机器字都会大大增加内存占用：

\begin{Verbatim}[commandchars=\\\{\},breaklines=true,breakanywhere=true,fontsize=\codefont,xleftmargin=17bp]
\PYG{k}{mod}\PYG{+w}{ }\PYG{n+nn}{ref\PYGZus{}with\PYGZus{}flag}\PYG{+w}{ }\PYG{p}{\PYGZob{}}
\PYG{+w}{    }\PYG{k}{use}\PYG{+w}{ }\PYG{n}{std}\PYG{p}{::}\PYG{n}{marker}\PYG{p}{::}\PYG{n}{PhantomData}\PYG{p}{;}
\PYG{+w}{    }\PYG{k}{use}\PYG{+w}{ }\PYG{n}{std}\PYG{p}{::}\PYG{n}{mem}\PYG{p}{::}\PYG{n}{align\PYGZus{}of}\PYG{p}{;}

\PYG{+w}{    }\PYG{l+s+sd}{/// 包装在单个机器字中的`\PYGZam{}T`和`bool`}
\PYG{+w}{    }\PYG{l+s+sd}{/// 类型`T`要求必须至少按两字节对齐}
\PYG{+w}{    }\PYG{l+s+sd}{///}
\PYG{+w}{    }\PYG{l+s+sd}{/// 如果你是那种中规中矩的程序员，从未想过还能从某个指针中偷出}
\PYG{+w}{    }\PYG{l+s+sd}{/// 第 20 位（数据的最低位），那么现在可以安全地做到这一点了！}
\PYG{+w}{    }\PYG{l+s+sd}{/// （“但这样做并不像想象中那么刺激啊……”）}
\PYG{+w}{    }\PYG{k}{pub}\PYG{+w}{ }\PYG{k}{struct}\PYG{+w}{ }\PYG{n+nc}{RefWithFlag}\PYG{o}{\PYGZlt{}}\PYG{o}{\PYGZsq{}}\PYG{n+na}{a}\PYG{p}{,}\PYG{+w}{ }\PYG{n}{T}\PYG{o}{\PYGZgt{}}\PYG{+w}{ }\PYG{p}{\PYGZob{}}
\PYG{+w}{        }\PYG{n}{ptr\PYGZus{}and\PYGZus{}bit}\PYG{p}{:}\PYG{+w}{ }\PYG{k+kt}{usize}\PYG{p}{,}
\PYG{+w}{        }\PYG{n}{behaves\PYGZus{}like}\PYG{p}{:}\PYG{+w}{ }\PYG{n+nc}{PhantomData}\PYG{o}{\PYGZlt{}}\PYG{o}{\PYGZam{}}\PYG{o}{\PYGZsq{}}\PYG{n+na}{a}\PYG{+w}{ }\PYG{n}{T}\PYG{o}{\PYGZgt{}}\PYG{+w}{ }\PYG{c+c1}{// 不占空间}
\PYG{+w}{    }\PYG{p}{\PYGZcb{}}

\PYG{+w}{    }\PYG{k}{impl}\PYG{o}{\PYGZlt{}}\PYG{o}{\PYGZsq{}}\PYG{n+na}{a}\PYG{p}{,}\PYG{+w}{ }\PYG{n}{T}\PYG{p}{:}\PYG{+w}{ }\PYG{o}{\PYGZsq{}}\PYG{n+na}{a}\PYG{o}{\PYGZgt{}}\PYG{+w}{ }\PYG{n}{RefWithFlag}\PYG{o}{\PYGZlt{}}\PYG{o}{\PYGZsq{}}\PYG{n+na}{a}\PYG{p}{,}\PYG{+w}{ }\PYG{n}{T}\PYG{o}{\PYGZgt{}}\PYG{+w}{ }\PYG{p}{\PYGZob{}}
\PYG{+w}{        }\PYG{k}{pub}\PYG{+w}{ }\PYG{k}{fn}\PYG{+w}{ }\PYG{n+nf}{new}\PYG{p}{(}\PYG{n}{ptr}\PYG{p}{:}\PYG{+w}{ }\PYG{k+kp}{\PYGZam{}}\PYG{o}{\PYGZsq{}}\PYG{n+na}{a}\PYG{+w}{ }\PYG{n+nc}{T}\PYG{p}{,}\PYG{+w}{ }\PYG{n}{flag}\PYG{p}{:}\PYG{+w}{ }\PYG{k+kt}{bool}\PYG{p}{)}\PYG{+w}{ }\PYG{p}{\PYGZhy{}\PYGZgt{}}\PYG{+w}{ }\PYG{n+nc}{RefWithFlag}\PYG{o}{\PYGZlt{}}\PYG{n}{T}\PYG{o}{\PYGZgt{}}\PYG{+w}{ }\PYG{p}{\PYGZob{}}
\PYG{+w}{            }\PYG{n+nf+fm}{assert!}\PYG{p}{(}\PYG{n}{align\PYGZus{}of}\PYG{p}{:}\PYG{p}{:}\PYG{o}{\PYGZlt{}}\PYG{n}{T}\PYG{o}{\PYGZgt{}}\PYG{p}{(}\PYG{p}{)}\PYG{+w}{ }\PYG{o}{\PYGZpc{}}\PYG{+w}{ }\PYG{l+m+mi}{2}\PYG{+w}{ }\PYG{o}{=}\PYG{o}{=}\PYG{+w}{ }\PYG{l+m+mi}{0}\PYG{p}{)}\PYG{p}{;}
\PYG{+w}{            }\PYG{n}{RefWithFlag}\PYG{+w}{ }\PYG{p}{\PYGZob{}}
\PYG{+w}{                }\PYG{n}{ptr\PYGZus{}and\PYGZus{}bit}\PYG{p}{:}\PYG{+w}{ }\PYG{n+nc}{ptr}\PYG{+w}{ }\PYG{k}{as}\PYG{+w}{ }\PYG{o}{*}\PYG{k}{const}\PYG{+w}{ }\PYG{n}{T}\PYG{+w}{ }\PYG{k}{as}\PYG{+w}{ }\PYG{k+kt}{usize}\PYG{+w}{ }\PYG{o}{|}\PYG{+w}{ }\PYG{n}{flag}\PYG{+w}{ }\PYG{k}{as}\PYG{+w}{ }\PYG{k+kt}{usize}\PYG{p}{,}
\PYG{+w}{                }\PYG{n}{behaves\PYGZus{}like}\PYG{p}{:}\PYG{+w}{ }\PYG{n+nc}{PhantomData}
\PYG{+w}{            }\PYG{p}{\PYGZcb{}}
\PYG{+w}{        }\PYG{p}{\PYGZcb{}}

\PYG{+w}{        }\PYG{k}{pub}\PYG{+w}{ }\PYG{k}{fn}\PYG{+w}{ }\PYG{n+nf}{get\PYGZus{}ref}\PYG{p}{(}\PYG{o}{\PYGZam{}}\PYG{n+nb+bp}{self}\PYG{p}{)}\PYG{+w}{ }\PYG{p}{\PYGZhy{}\PYGZgt{}}\PYG{+w}{ }\PYG{k+kp}{\PYGZam{}}\PYG{o}{\PYGZsq{}}\PYG{n+na}{a}\PYG{+w}{ }\PYG{n+nc}{T}\PYG{+w}{ }\PYG{p}{\PYGZob{}}
\PYG{+w}{            }\PYG{k}{unsafe}\PYG{+w}{ }\PYG{p}{\PYGZob{}}
\PYG{+w}{                }\PYG{k+kd}{let}\PYG{+w}{ }\PYG{n}{ptr}\PYG{+w}{ }\PYG{o}{=}\PYG{+w}{ }\PYG{p}{(}\PYG{n+nb+bp}{self}\PYG{p}{.}\PYG{n}{ptr\PYGZus{}and\PYGZus{}bit}\PYG{+w}{ }\PYG{o}{\PYGZam{}}\PYG{+w}{ }\PYG{o}{!}\PYG{l+m+mi}{1}\PYG{p}{)}\PYG{+w}{ }\PYG{k}{as}\PYG{+w}{ }\PYG{o}{*}\PYG{k}{const}\PYG{+w}{ }\PYG{n}{T}\PYG{p}{;}
\PYG{+w}{                }\PYG{o}{\PYGZam{}}\PYG{o}{*}\PYG{n}{ptr}
\PYG{+w}{            }\PYG{p}{\PYGZcb{}}
\PYG{+w}{        }\PYG{p}{\PYGZcb{}}

\PYG{+w}{        }\PYG{k}{pub}\PYG{+w}{ }\PYG{k}{fn}\PYG{+w}{ }\PYG{n+nf}{get\PYGZus{}flag}\PYG{p}{(}\PYG{o}{\PYGZam{}}\PYG{n+nb+bp}{self}\PYG{p}{)}\PYG{+w}{ }\PYG{p}{\PYGZhy{}\PYGZgt{}}\PYG{+w}{ }\PYG{k+kt}{bool}\PYG{+w}{ }\PYG{p}{\PYGZob{}}
\PYG{+w}{            }\PYG{n+nb+bp}{self}\PYG{p}{.}\PYG{n}{ptr\PYGZus{}and\PYGZus{}bit}\PYG{+w}{ }\PYG{o}{\PYGZam{}}\PYG{+w}{ }\PYG{l+m+mi}{1}\PYG{+w}{ }\PYG{o}{!}\PYG{o}{=}\PYG{+w}{ }\PYG{l+m+mi}{0}
\PYG{+w}{        }\PYG{p}{\PYGZcb{}}
\PYG{+w}{    }\PYG{p}{\PYGZcb{}}
\PYG{p}{\PYGZcb{}}
\end{Verbatim}

这段代码利用了这样一个事实，即许多类型在内存中必须放置在偶数地址：由于偶数地址的最低有效位始终为 0，因此可以在那里存储其他内容，然后通过屏蔽最低位来可靠地重建原始地址。并非所有类型都符合条件，比如类型 \texttt{u8} 和 \texttt{(bool, [i8; 2])} 可以放在任何地址。但是我们可以检查此类型在构造方面的对齐情况，并拒绝不适用的类型。

可以像下面这样使用 \texttt{RefWithFlag}：

\begin{Verbatim}[commandchars=\\\{\},breaklines=true,breakanywhere=true,fontsize=\codefont,xleftmargin=17bp]
\PYG{k}{use}\PYG{+w}{ }\PYG{n}{ref\PYGZus{}with\PYGZus{}flag}\PYG{p}{::}\PYG{n}{RefWithFlag}\PYG{p}{;}

\PYG{k+kd}{let}\PYG{+w}{ }\PYG{n}{vec}\PYG{+w}{ }\PYG{o}{=}\PYG{+w}{ }\PYG{n+nf+fm}{vec!}\PYG{p}{[}\PYG{l+m+mi}{10}\PYG{p}{,}\PYG{+w}{ }\PYG{l+m+mi}{20}\PYG{p}{,}\PYG{+w}{ }\PYG{l+m+mi}{30}\PYG{p}{]}\PYG{p}{;}
\PYG{k+kd}{let}\PYG{+w}{ }\PYG{n}{flagged}\PYG{+w}{ }\PYG{o}{=}\PYG{+w}{ }\PYG{n}{RefWithFlag}\PYG{p}{::}\PYG{n}{new}\PYG{p}{(}\PYG{o}{\PYGZam{}}\PYG{n}{vec}\PYG{p}{,}\PYG{+w}{ }\PYG{k+kc}{true}\PYG{p}{)}\PYG{p}{;}
\PYG{n+nf+fm}{assert\PYGZus{}eq!}\PYG{p}{(}\PYG{n}{flagged}\PYG{p}{.}\PYG{n}{get\PYGZus{}ref}\PYG{p}{(}\PYG{p}{)}\PYG{p}{[}\PYG{l+m+mi}{1}\PYG{p}{]}\PYG{p}{,}\PYG{+w}{ }\PYG{l+m+mi}{20}\PYG{p}{)}\PYG{p}{;}
\PYG{n+nf+fm}{assert\PYGZus{}eq!}\PYG{p}{(}\PYG{n}{flagged}\PYG{p}{.}\PYG{n}{get\PYGZus{}flag}\PYG{p}{(}\PYG{p}{)}\PYG{p}{,}\PYG{+w}{ }\PYG{k+kc}{true}\PYG{p}{)}\PYG{p}{;}
\end{Verbatim}

构造函数 \texttt{RefWithFlag::new} 会接受一个引用和一个 \texttt{bool} 值，并断言此引用具有适当的类型，然后把它转换为裸指针，再转换为 \texttt{usize} 类型。\texttt{usize} 类型大小的定义是足够在我们正在编译的任何处理器上保存一个指针，因此将裸指针转换为 \texttt{usize} 并返回它是有明确定义的。一旦有了 \texttt{usize}，我们就知道它必然是偶数，所以可以使用按位或运算符 \texttt{|} 将其与已转换为整数 0 或 1 的 \texttt{bool} 值组合起来。

\texttt{get\_flag} 方法用于提取 \texttt{RefWithFlag} 的 \texttt{bool} 部分。这很简单：只要取出最低位并检查结果是否非零就可以了（\texttt{self.ptr\_and\_bit \& 1 != 0}）。

\texttt{get\_ref} 方法用于从 \texttt{RefWithFlag} 中提取引用。首先，它会屏蔽 \texttt{usize} 的最低位（\texttt{self.ptr\_and\_bit \& !1}）并将其转换为裸指针。\texttt{as} 运算符无法将裸指针转换为引用，但我们可以解引用裸指针（当然是在 \texttt{unsafe} 块中）并借用它。借用一个裸指针的引用目标会得到一个无限生命周期的引用：Rust 会赋予引用任何生命周期来检查它周围的代码（如果有的话）。但是，通常还有一些更准确的特定生命周期，因此会发现更多错误。在这个例子中，由于 \texttt{get\_ref} 的返回类型是 \texttt{\&'a T}，因此 Rust 认为该引用的生命周期与 \texttt{RefWithFlag} 的生命周期参数 \texttt{'a} 相同，这正是我们想要的，因为这个生命周期就是最初那个引用的生命周期。

在内存中，\texttt{RefWithFlag} 看起来很像 \texttt{usize}：由于 \texttt{PhantomData}（意思是虚构的数据）是零大小的类型，因此 \texttt{behaves\_like} 字段并不会占用结构体中的空间。但是，为了让 Rust 知道该如何处理使用 \texttt{RefWithFlag} 的代码中的生命周期，\texttt{PhantomData} 是必需的。想象一下没有 \texttt{behaves\_like} 字段的类型会是什么样子：

\begin{Verbatim}[commandchars=\\\{\},breaklines=true,breakanywhere=true,fontsize=\codefont,xleftmargin=17bp]
\PYG{c+c1}{// 这无法编译}
\PYG{k}{pub}\PYG{+w}{ }\PYG{k}{struct}\PYG{+w}{ }\PYG{n+nc}{RefWithFlag}\PYG{o}{\PYGZlt{}}\PYG{o}{\PYGZsq{}}\PYG{n+na}{a}\PYG{p}{,}\PYG{+w}{ }\PYG{n}{T}\PYG{p}{:}\PYG{+w}{ }\PYG{o}{\PYGZsq{}}\PYG{n+na}{a}\PYG{o}{\PYGZgt{}}\PYG{+w}{ }\PYG{p}{\PYGZob{}}
\PYG{+w}{    }\PYG{n}{ptr\PYGZus{}and\PYGZus{}bit}\PYG{p}{:}\PYG{+w}{ }\PYG{k+kt}{usize}
\PYG{p}{\PYGZcb{}}
\end{Verbatim}

如第 5 章所述，任何包含引用的结构体，其生命周期都不能超出它们借用的值，以免引用变成悬空指针。这个结构体必须遵守适用于其字段的限制。这当然也适用于 \texttt{RefWithFlag}：在刚刚看到的示例代码中，\texttt{flagged} 的生命周期不能超出 \texttt{vec}，因为 \texttt{flagged.get\_ref()} 会返回对它的引用。但是我们简化版的 \texttt{RefWithFlag} 类型根本不包含任何引用，并且从不使用其生命周期参数 \texttt{'a}，因为这只是一个 \texttt{usize}。怎么让 Rust 知道应该如何限制 \texttt{flagged} 的生命周期呢？包含一个 \texttt{PhantomData<\&'a T>} 字段就是为了告诉 Rust 应该将 \texttt{RefWithFlag<'a, T>} \textbf{视为}包含一个 \texttt{\&'a T}，却不会实际影响此结构体的表示方式。

尽管 Rust 并不真正知道发生了什么（这就是 \texttt{RefWithFlag} 不安全的原因），但它会尽力帮助你解决这个问题。如果省略了 \texttt{behaves\_like} 字段，那么 Rust 就会报错说参数 \texttt{'a} 和 \texttt{T} 未使用，并建议使用 \texttt{PhantomData}。

\texttt{RefWithFlag} 使用了与之前介绍的 \texttt{Ascii} 类型相同的策略来避免其 \texttt{unsafe} 块中的未定义行为。类型本身是 \texttt{pub} 的，但其字段不是，这意味着只有 \texttt{ref\_with\_flag} 模块中的代码才能创建或查看 \texttt{RefWithFlag} 值。你不必检查太多代码就可以确信 \texttt{ptr\_and\_bit} 字段是构造良好的。

\subsection{可空指针}

Rust 中的空裸指针是一个零地址，与 C 和 C++ 中一样。对于任意类型 \texttt{T}，\texttt{std::ptr::null<T>} 函数会返回一个 \texttt{*const T} 空指针，而 \texttt{std::ptr::null\_mut<T>} 会返回一个 \texttt{*mut T} 空指针。

检查裸指针是否为空有几种方法。最简单的是 \texttt{is\_null} 方法，但 \texttt{as\_ref} 方法可能更方便。\texttt{as\_ref} 方法会接受 \texttt{*const T} 指针并返回 \texttt{Option<\&'a T>}，以便将一个空指针变成 \texttt{None}。同样，\texttt{as\_mut} 方法会将 \texttt{*mut T} 指针转换为 \texttt{Option<\&'a mut T>} 值。

\subsection{类型大小与对齐方式}

任何固定大小类型（\texttt{Sized}）的值都会在内存中占用固定数量的字节，并且必须放置在由机器体系结构决定的某个\textbf{对齐}值的倍数的地址处。例如，一个 \texttt{(i32, i32)} 元组占用 8 字节，而大多数处理器更喜欢将其放置在 4 的倍数地址处。

调用 \texttt{std::mem::size\_of::<T>()} 会返回类型 \texttt{T} 值的大小（以字节为单位），而调用 \texttt{std::mem::align\_of::<T>()} 会返回其所需的对齐方式。例如：

\begin{Verbatim}[commandchars=\\\{\},breaklines=true,breakanywhere=true,fontsize=\codefont,xleftmargin=17bp]
\PYG{n+nf+fm}{assert\PYGZus{}eq!}\PYG{p}{(}\PYG{n}{std}\PYG{p}{::}\PYG{n}{mem}\PYG{p}{::}\PYG{n}{size\PYGZus{}of}\PYG{p}{:}\PYG{p}{:}\PYG{o}{\PYGZlt{}}\PYG{k+kt}{i64}\PYG{o}{\PYGZgt{}}\PYG{p}{(}\PYG{p}{)}\PYG{p}{,}\PYG{+w}{ }\PYG{l+m+mi}{8}\PYG{p}{)}\PYG{p}{;}
\PYG{n+nf+fm}{assert\PYGZus{}eq!}\PYG{p}{(}\PYG{n}{std}\PYG{p}{::}\PYG{n}{mem}\PYG{p}{::}\PYG{n}{align\PYGZus{}of}\PYG{p}{:}\PYG{p}{:}\PYG{o}{\PYGZlt{}}\PYG{p}{(}\PYG{k+kt}{i32}\PYG{p}{,}\PYG{+w}{ }\PYG{k+kt}{i32}\PYG{p}{)}\PYG{o}{\PYGZgt{}}\PYG{p}{(}\PYG{p}{)}\PYG{p}{,}\PYG{+w}{ }\PYG{l+m+mi}{4}\PYG{p}{)}\PYG{p}{;}
\end{Verbatim}

任何类型总是对齐到二的 \emph{n} 次幂。

即使在技术上可以填入更小的空间，类型的大小也总是会四舍五入为其对齐方式的倍数。例如，尽管像 \texttt{(f32, u8)} 这样的元组只需要 5 字节，但 \texttt{size\_of::<(f32, u8)>()} 是 \texttt{8}，因为 \texttt{align\_of::<(f32, u8)>()} 是 \texttt{4}。这会确保如果你有一个数组，那么元素类型的大小总能反映出一个元素与其下一个元素的间距。

对于无固定大小类型，其大小和对齐方式取决于手头的值。给定对无固定大小值的引用，\texttt{std::mem::size\_of\_val} 函数和 \texttt{std::mem::align\_of\_val} 函数会返回值的大小和对齐方式。这两个函数可以对固定大小类型和无固定大小类型的引用进行操作。

\begin{Verbatim}[commandchars=\\\{\},breaklines=true,breakanywhere=true,fontsize=\codefont,xleftmargin=17bp]
\PYG{c+c1}{// 指向切片的胖指针包含其引用目标的长度}
\PYG{k+kd}{let}\PYG{+w}{ }\PYG{n}{slice}\PYG{p}{:}\PYG{+w}{ }\PYG{k+kp}{\PYGZam{}}\PYG{p}{[}\PYG{k+kt}{i32}\PYG{p}{]}\PYG{+w}{ }\PYG{o}{=}\PYG{+w}{ }\PYG{o}{\PYGZam{}}\PYG{p}{[}\PYG{l+m+mi}{1}\PYG{p}{,}\PYG{+w}{ }\PYG{l+m+mi}{3}\PYG{p}{,}\PYG{+w}{ }\PYG{l+m+mi}{9}\PYG{p}{,}\PYG{+w}{ }\PYG{l+m+mi}{27}\PYG{p}{,}\PYG{+w}{ }\PYG{l+m+mi}{81}\PYG{p}{]}\PYG{p}{;}
\PYG{n+nf+fm}{assert\PYGZus{}eq!}\PYG{p}{(}\PYG{n}{std}\PYG{p}{::}\PYG{n}{mem}\PYG{p}{::}\PYG{n}{size\PYGZus{}of\PYGZus{}val}\PYG{p}{(}\PYG{n}{slice}\PYG{p}{)}\PYG{p}{,}\PYG{+w}{ }\PYG{l+m+mi}{20}\PYG{p}{)}\PYG{p}{;}

\PYG{k+kd}{let}\PYG{+w}{ }\PYG{n}{text}\PYG{p}{:}\PYG{+w}{ }\PYG{k+kp}{\PYGZam{}}\PYG{k+kt}{str}\PYG{+w}{ }\PYG{o}{=}\PYG{+w}{ }\PYG{l+s}{\PYGZdq{}}\PYG{l+s}{alligator}\PYG{l+s}{\PYGZdq{}}\PYG{p}{;}
\PYG{n+nf+fm}{assert\PYGZus{}eq!}\PYG{p}{(}\PYG{n}{std}\PYG{p}{::}\PYG{n}{mem}\PYG{p}{::}\PYG{n}{size\PYGZus{}of\PYGZus{}val}\PYG{p}{(}\PYG{n}{text}\PYG{p}{)}\PYG{p}{,}\PYG{+w}{ }\PYG{l+m+mi}{9}\PYG{p}{)}\PYG{p}{;}

\PYG{k}{use}\PYG{+w}{ }\PYG{n}{std}\PYG{p}{::}\PYG{n}{fmt}\PYG{p}{::}\PYG{n}{Display}\PYG{p}{;}
\PYG{k+kd}{let}\PYG{+w}{ }\PYG{n}{unremarkable}\PYG{p}{:}\PYG{+w}{ }\PYG{k+kp}{\PYGZam{}}\PYG{n+nc}{dyn}\PYG{+w}{ }\PYG{n}{Display}\PYG{+w}{ }\PYG{o}{=}\PYG{+w}{ }\PYG{o}{\PYGZam{}}\PYG{l+m+mi}{193\PYGZus{}}\PYG{k}{u8}\PYG{p}{;}
\PYG{k+kd}{let}\PYG{+w}{ }\PYG{n}{remarkable}\PYG{p}{:}\PYG{+w}{ }\PYG{k+kp}{\PYGZam{}}\PYG{n+nc}{dyn}\PYG{+w}{ }\PYG{n}{Display}\PYG{+w}{ }\PYG{o}{=}\PYG{+w}{ }\PYG{o}{\PYGZam{}}\PYG{l+m+mf}{0.0072973525664}\PYG{p}{;}

\PYG{c+c1}{// 这些会返回特型对象指向的值的大小/对齐方式，而不是特型对象}
\PYG{c+c1}{// 本身的大小/对齐方式。此信息来自特型对象引用的虚表}
\PYG{n+nf+fm}{assert\PYGZus{}eq!}\PYG{p}{(}\PYG{n}{std}\PYG{p}{::}\PYG{n}{mem}\PYG{p}{::}\PYG{n}{size\PYGZus{}of\PYGZus{}val}\PYG{p}{(}\PYG{n}{unremarkable}\PYG{p}{)}\PYG{p}{,}\PYG{+w}{ }\PYG{l+m+mi}{1}\PYG{p}{)}\PYG{p}{;}
\PYG{n+nf+fm}{assert\PYGZus{}eq!}\PYG{p}{(}\PYG{n}{std}\PYG{p}{::}\PYG{n}{mem}\PYG{p}{::}\PYG{n}{align\PYGZus{}of\PYGZus{}val}\PYG{p}{(}\PYG{n}{remarkable}\PYG{p}{)}\PYG{p}{,}\PYG{+w}{ }\PYG{l+m+mi}{8}\PYG{p}{)}\PYG{p}{;}
\end{Verbatim}

\subsection{指针运算}

Rust 会将数组、切片或向量的元素排布为单个连续的内存块，如图 22-1 所示。元素的间隔很均匀，因此如果每个元素占用 \texttt{size} 字节，则第 \texttt{i} 个元素就从第 \texttt{i * size} 字节开始。

\myFigure{images/ch22/fig-22-1.jpeg}{图 22-1：内存中的数组}

这样做有一个好处：如果你有两个指向数组元素的裸指针，那么比较指针就会得到与比较元素索引相同的结果。如果 \texttt{i < j}，则指向第 \texttt{i} 个元素的裸指针一定小于指向第 \texttt{j} 个元素的裸指针。这使得裸指针可用作数组遍历的边界。事实上，标准库对切片的简单迭代器最初就是这样定义的：

\begin{Verbatim}[commandchars=\\\{\},breaklines=true,breakanywhere=true,fontsize=\codefont,xleftmargin=17bp]
\PYG{k}{struct}\PYG{+w}{ }\PYG{n+nc}{Iter}\PYG{o}{\PYGZlt{}}\PYG{o}{\PYGZsq{}}\PYG{n+na}{a}\PYG{p}{,}\PYG{+w}{ }\PYG{n}{T}\PYG{o}{\PYGZgt{}}\PYG{+w}{ }\PYG{p}{\PYGZob{}}
\PYG{+w}{    }\PYG{n}{ptr}\PYG{p}{:}\PYG{+w}{ }\PYG{o}{*}\PYG{k}{const}\PYG{+w}{ }\PYG{n}{T}\PYG{p}{,}
\PYG{+w}{    }\PYG{n}{end}\PYG{p}{:}\PYG{+w}{ }\PYG{o}{*}\PYG{k}{const}\PYG{+w}{ }\PYG{n}{T}\PYG{p}{,}
\PYG{+w}{    }\PYG{o}{..}\PYG{p}{.}
\PYG{p}{\PYGZcb{}}
\end{Verbatim}

\texttt{ptr} 字段指向迭代应该生成的下一个元素，\texttt{end} 字段作为界限：当 \texttt{ptr == end} 时，迭代完成。

数组布局的另一个好处是：如果 \texttt{element\_ptr} 是指向某个数组的第 \texttt{i} 个元素的 \texttt{*const T} 或 \texttt{*mut T} 裸指针，那么 \texttt{element\_ptr.offset(o)} 就是指向第 \texttt{(i + o)} 个元素的裸指针。它的定义等效于如下内容：

\begin{Verbatim}[commandchars=\\\{\},breaklines=true,breakanywhere=true,fontsize=\codefont,xleftmargin=17bp]
\PYG{k}{fn}\PYG{+w}{ }\PYG{n+nf}{offset}\PYG{o}{\PYGZlt{}}\PYG{n}{T}\PYG{o}{\PYGZgt{}}\PYG{p}{(}\PYG{n}{ptr}\PYG{p}{:}\PYG{+w}{ }\PYG{o}{*}\PYG{k}{const}\PYG{+w}{ }\PYG{n}{T}\PYG{p}{,}\PYG{+w}{ }\PYG{n}{count}\PYG{p}{:}\PYG{+w}{ }\PYG{k+kt}{isize}\PYG{p}{)}\PYG{+w}{ }\PYG{p}{\PYGZhy{}\PYGZgt{}}\PYG{+w}{ }\PYG{o}{*}\PYG{k}{const}\PYG{+w}{ }\PYG{n}{T}
\PYG{+w}{    }\PYG{k}{where}\PYG{+w}{ }\PYG{n}{T}\PYG{p}{:}\PYG{+w}{ }\PYG{n+nb}{Sized}
\PYG{p}{\PYGZob{}}
\PYG{+w}{    }\PYG{k+kd}{let}\PYG{+w}{ }\PYG{n}{bytes\PYGZus{}per\PYGZus{}element}\PYG{+w}{ }\PYG{o}{=}\PYG{+w}{ }\PYG{n}{std}\PYG{p}{::}\PYG{n}{mem}\PYG{p}{::}\PYG{n}{size\PYGZus{}of}\PYG{p}{:}\PYG{p}{:}\PYG{o}{\PYGZlt{}}\PYG{n}{T}\PYG{o}{\PYGZgt{}}\PYG{p}{(}\PYG{p}{)}\PYG{+w}{ }\PYG{k}{as}\PYG{+w}{ }\PYG{k+kt}{isize}\PYG{p}{;}
\PYG{+w}{    }\PYG{k+kd}{let}\PYG{+w}{ }\PYG{n}{byte\PYGZus{}offset}\PYG{+w}{ }\PYG{o}{=}\PYG{+w}{ }\PYG{n}{count}\PYG{+w}{ }\PYG{o}{*}\PYG{+w}{ }\PYG{n}{bytes\PYGZus{}per\PYGZus{}element}\PYG{p}{;}
\PYG{+w}{    }\PYG{p}{(}\PYG{n}{ptr}\PYG{+w}{ }\PYG{k}{as}\PYG{+w}{ }\PYG{k+kt}{isize}\PYG{p}{)}\PYG{p}{.}\PYG{n}{checked\PYGZus{}add}\PYG{p}{(}\PYG{n}{byte\PYGZus{}offset}\PYG{p}{)}\PYG{p}{.}\PYG{n}{unwrap}\PYG{p}{(}\PYG{p}{)}\PYG{+w}{ }\PYG{k}{as}\PYG{+w}{ }\PYG{o}{*}\PYG{k}{const}\PYG{+w}{ }\PYG{n}{T}
\PYG{p}{\PYGZcb{}}
\end{Verbatim}

\texttt{std::mem::size\_of::<T>} 函数会返回类型 \texttt{T} 的字节大小。根据定义，由于 \texttt{isize} 大到足以容纳一个地址，因此可以将基指针转换为 \texttt{isize}，对得到的值进行算术运算，然后将结果转换回指针。

可以生成指向数组末尾之后第一字节的指针。虽然不能对这样的指针解引用，但可以用它来表示循环的界限或用于边界检查。

但是，使用 \texttt{offset} 生成超出该点或指向数组开头之前的指针是未定义行为，即使从未对它解引用也是如此。为了方便优化，Rust 会假设当 \texttt{i} 为正值时 \texttt{ptr.offset(i) > ptr}，当 \texttt{i} 为负值时 \texttt{ptr.offset(i) < ptr}。这个假设似乎是安全的，但如果 \texttt{offset} 中的算术溢出了 \texttt{isize} 值，那么可能就不成立了。如果把 \texttt{i} 限制在 \texttt{ptr} 的同一个数组范围内，则肯定不会发生溢出：毕竟数组本身不会溢出地址空间的边界。（为了让指向结尾之后第一字节的指针安全，Rust 从来都不会将值放在地址空间的上端。）

如果确实需要将指针偏移到与其关联的数组的界限之外，则可以使用 \texttt{wrapping\_offset} 方法。该方法与 \texttt{offset} 等效，但 Rust 不会假设 \texttt{ptr.wrapping\_offset(i)} 和 \texttt{ptr} 本身的相对顺序。当然，你仍然不能对此类指针解引用，除非确信它们会落在数组中。

\subsection{移动入和移动出内存}

如果你正在实现的类型需要管理自己的内存，那么就要跟踪内存中哪些部分保存了有效值，而哪些是未初始化的，就像 Rust 处理局部变量一样。考虑下面这段代码：

\begin{Verbatim}[commandchars=\\\{\},breaklines=true,breakanywhere=true,fontsize=\codefont,xleftmargin=17bp]
\PYG{k+kd}{let}\PYG{+w}{ }\PYG{n}{pot}\PYG{+w}{ }\PYG{o}{=}\PYG{+w}{ }\PYG{l+s}{\PYGZdq{}}\PYG{l+s}{pasta}\PYG{l+s}{\PYGZdq{}}\PYG{p}{.}\PYG{n}{to\PYGZus{}string}\PYG{p}{(}\PYG{p}{)}\PYG{p}{;}
\PYG{k+kd}{let}\PYG{+w}{ }\PYG{n}{plate}\PYG{+w}{ }\PYG{o}{=}\PYG{+w}{ }\PYG{n}{pot}\PYG{p}{;}
\end{Verbatim}

上述代码运行后，情况如图 22-2 所示。

\myFigure{images/ch22/fig-22-2.jpeg}{图 22-2：将字符串从一个局部变量转移给另一个局部变量}

赋值后，\texttt{pot} 处于未初始化状态，而 \texttt{plate} 成了字符串的拥有者。

在机器层面，没有指定移动对源值的作用，但实际上它通常什么都不做。该赋值可能会使 \texttt{pot} 仍然保存着字符串的指针、容量和长度。当然，如果继续将其视为有效值将是灾难性的，但 Rust 会确保你不会这样做。

同样的考虑也适用于管理自己内存的数据结构。假设你运行了下面这段代码：

\begin{Verbatim}[commandchars=\\\{\},breaklines=true,breakanywhere=true,fontsize=\codefont,xleftmargin=17bp]
\PYG{k+kd}{let}\PYG{+w}{ }\PYG{k}{mut}\PYG{+w}{ }\PYG{n}{noodles}\PYG{+w}{ }\PYG{o}{=}\PYG{+w}{ }\PYG{n+nf+fm}{vec!}\PYG{p}{[}\PYG{l+s}{\PYGZdq{}}\PYG{l+s}{udon}\PYG{l+s}{\PYGZdq{}}\PYG{p}{.}\PYG{n}{to\PYGZus{}string}\PYG{p}{(}\PYG{p}{)}\PYG{p}{]}\PYG{p}{;}
\PYG{k+kd}{let}\PYG{+w}{ }\PYG{n}{soba}\PYG{+w}{ }\PYG{o}{=}\PYG{+w}{ }\PYG{l+s}{\PYGZdq{}}\PYG{l+s}{soba}\PYG{l+s}{\PYGZdq{}}\PYG{p}{.}\PYG{n}{to\PYGZus{}string}\PYG{p}{(}\PYG{p}{)}\PYG{p}{;}
\PYG{k+kd}{let}\PYG{+w}{ }\PYG{n}{last}\PYG{p}{;}
\end{Verbatim}

在内存中，状态如图 22-3 所示。

\myFigure{images/ch22/fig-22-3.jpeg}{图 22-3：具有未初始化的空闲容量的向量}

这个向量有空闲容量可以再容纳一个元素，但空闲容量中存放的是垃圾数据，可能是以前的内存残余。假设你随后运行了如下代码：

\begin{Verbatim}[commandchars=\\\{\},breaklines=true,breakanywhere=true,fontsize=\codefont,xleftmargin=17bp]
\PYG{n}{noodles}\PYG{p}{.}\PYG{n}{push}\PYG{p}{(}\PYG{n}{soba}\PYG{p}{)}\PYG{p}{;}
\end{Verbatim}

将字符串压入向量会将未初始化的内存转换为新元素，如图 22-4 所示。

\myFigure{images/ch22/fig-22-4.jpeg}{图 22-4：将 \texttt{soba} 的值推入向量之后}

该向量已初始化其空白空间，以便拥有该字符串，并增加其长度，以便将其标记为新的有效元素。向量现在是字符串的拥有者，你可以引用它的第二个元素了，而丢弃此向量将释放两个字符串。\texttt{soba} 现在处于未初始化状态。

最后，考虑一下当从向量中弹出一个值时会发生什么：

\begin{Verbatim}[commandchars=\\\{\},breaklines=true,breakanywhere=true,fontsize=\codefont,xleftmargin=17bp]
\PYG{+w}{    }\PYG{n}{last}\PYG{+w}{ }\PYG{o}{=}\PYG{+w}{ }\PYG{n}{noodles}\PYG{p}{.}\PYG{n}{pop}\PYG{p}{(}\PYG{p}{)}\PYG{p}{.}\PYG{n}{unwrap}\PYG{p}{(}\PYG{p}{)}\PYG{p}{;}
\end{Verbatim}

在内存中，现在看起来如图 22-5 所示。

\myFigure{images/ch22/fig-22-5.jpeg}{图 22-5：把向量中的一个元素弹出到 \texttt{last} 之后}

变量 \texttt{last} 取得了字符串的所有权。向量已减小其 \texttt{length} 以指示用于保存字符串的空间现在未初始化。

就像之前的 \texttt{pot} 和 \texttt{pasta} 一样，\texttt{soba}、\texttt{last} 和向量的可用空间这三者可能存有相同的位模式。但只有 \texttt{last} 被认为拥有这个值。将其他两个位置中的任何一个视为有效位置都是错误的。

初始化值的真正定义是\textbf{应视为有效}的值。写入值的字节通常是初始化的必要部分，但这只是为了将其视为有效的而做的准备工作。移动和复制对内存的影响是一样的，两者之间的区别在于，在移动之后，源不再被视为有效值，而在复制之后，源和目标都处于有效状态。

Rust 会在编译期跟踪哪些局部变量处于有效状态，并阻止你使用值已转移给其他地方的变量。\texttt{Vec}、\texttt{HashMap}、\texttt{Box} 等类型会动态跟踪它们的缓冲区。如果你实现了一个管理自己内存的类型，则也需要这样做。

Rust 为实现这些类型提供了两个基本操作。

\begin{mySidenote}
\texttt{std::ptr::read(src)}（读取）
\end{mySidenote}

将值移动出 \texttt{src} 指向的位置，将所有权转移给调用者。\texttt{src} 参数应该是一个 \texttt{*const T} 裸指针，其中 \texttt{T} 是一个固定大小类型。调用此函数后，\texttt{*src} 的内容不受影响，但除非 \texttt{T} 是 \texttt{Copy} 类型，否则你必须确保自己的程序会将它们视为未初始化内存。

这是 \texttt{Vec::pop} 背后的操作。要弹出一个值，就要调用 \texttt{read} 将该值移出缓冲区，然后递减长度以将该空间标记为未初始化容量。

\begin{myNote}
\texttt{std::ptr::write(dest, value)}（写入）
\end{myNote}

将 \texttt{value} 转移给 \texttt{dest} 指向的位置，该位置在调用之前必须是未初始化内存。引用目标现在拥有该值。在这里，\texttt{dest} 必须是一个 \texttt{*mut T} 裸指针并且 \texttt{value} 是一个 \texttt{T} 值，其中 \texttt{T} 是固定大小类型。

这就是 \texttt{Vec::push} 背后的操作。压入一个值会调用 \texttt{write} 将值转移给下一个可用空间，然后增加长度以将该空间标记为有效元素。

两者都是自由函数，而不是裸指针类型的方法。

请注意，不能使用任何 Rust 的安全指针类型来执行这些操作。安全指针类型会要求其引用目标始终是初始化的，因此将未初始化内存转换为值或相反的操作都超出了它们的能力范围。而裸指针符合这种要求。

标准库还提供了将值数组从一个内存块移动到另一个内存块的函数。

\begin{myNote}
\texttt{std::ptr::copy(src, dst, count)}（复制）
\end{myNote}

将内存中从 \texttt{src} 开始的 \texttt{count} 个值数组移动到 \texttt{dst} 处，就像编写了一个 \texttt{read} 和 \texttt{write} 调用循环以一次性移动它们一样。调用之前目标内存必须是未初始化的，调用之后源内存要保持未初始化状态。\texttt{src} 参数和 \texttt{dest} 参数必须是 \texttt{*const T} 裸指针和 \texttt{*mut T} 裸指针，并且 \texttt{count} 必须是 \texttt{usize}。

\begin{myNote}
\texttt{ptr.copy\_to(dst, count)}（复制到）
\end{myNote}

一个更方便的 \texttt{copy} 版本，它会将内存中从 \texttt{ptr} 开始的 \texttt{count} 个值的数组转移给 \texttt{dst}，而不用以其起点作为参数。

\begin{myNote}
\texttt{std::ptr::copy\_nonoverlapping(src, dst, count)}（复制，无重叠版）
\end{myNote}

就像对 \texttt{copy} 的类似调用一样，但是它的契约进一步要求源内存块和目标内存块不能重叠。这可能比调用 \texttt{copy} 略微快一些。

\begin{myNote}
\texttt{ptr.copy\_to\_nonoverlapping(dst, count)}（复制到，无重叠版）
\end{myNote}

一个更方便的 \texttt{copy\_nonoverlapping} 版本，就像 \texttt{copy\_to}。

还有另外两组 \texttt{read} 函数和 \texttt{write} 函数，它们也位于 \texttt{std::ptr} 模块中。

\begin{myNote}
\texttt{read\_unaligned}（读取，未对齐版）和 \texttt{write\_unaligned}（写入，未对齐版）
\end{myNote}

与 \texttt{read} 和 \texttt{write} 类似，但是这两个函数的指针不需要像引用目标类型通常要求的那样对齐。它们可能比普通的 \texttt{read} 函数和 \texttt{write} 函数慢一点儿。

\begin{myNote}
\texttt{read\_volatile}（读取，易变版）和 \texttt{write\_volatile}（写入，易变版）
\end{myNote}

这两个函数对应于 C 或 C++ 中的易变（\texttt{volatile}）读取和易变写入。

\subsection{示例：\texttt{GapBuffer}}

下面是一个使用刚刚讲过的裸指针函数的示例。

假设你正在编写一个文本编辑器，并且正在寻找一种类型来表示文本。可以选择 \texttt{String} 并使用 \texttt{insert} 方法和 \texttt{remove} 方法在用户键入时插入字符和移除字符。但是如果在一个大文件的开头编辑文本，则这些方法可能开销会很高：插入新字符需要在内存中将整个字符串的其余部分都移到右侧，而删除则要将其全部移回左侧。你希望此类常见操作的开销低一些。

Emacs 文本编辑器使用了一种称为\textbf{间隙缓冲区}的简单数据结构，该数据结构可以在恒定时间内插入字符和删除字符。\texttt{String} 会将其所有空闲容量保留在文本的末尾，这使得 \texttt{push} 和 \texttt{pop} 的开销变得更低，而间隙缓冲区会将其空闲容量保留在文本中间，即正在进行编辑的位置。这种空闲容量称为\textbf{间隙}。在间隙处插入元素或删除元素的开销很低，只要根据需要缩小或扩大间隙即可。可以通过将文本从间隙的一侧移动到另一侧，来让间隙移动到你喜欢的任何位置。当间隙为空时，就迁移到更大的缓冲区。

虽然间隙缓冲区中的插入和删除速度很快，但如果想更改这些操作发生的位置就要将间隙移动到新位置。移动元素需要的时间与移动的距离成正比。幸运的是，典型的编辑活动通常都会在转移到别处之前，在缓冲区的临近区域中进行一系列更改。

本节将在 Rust 中实现间隙缓冲区。为了避免被 UTF-8 分散注意力，我们会让该缓冲区直接存储 \texttt{char} 值，但即使以其他形式存储文本，这些操作的原则也是一样的。

首先，我们会展示间隙缓冲区的实际应用。下列代码会创建一个 \texttt{GapBuffer}，在其中插入一些文本，然后将插入点移动到最后一个单词之前：

\begin{Verbatim}[commandchars=\\\{\},breaklines=true,breakanywhere=true,fontsize=\codefont,xleftmargin=17bp]
\PYG{k+kd}{let}\PYG{+w}{ }\PYG{k}{mut}\PYG{+w}{ }\PYG{n}{buf}\PYG{+w}{ }\PYG{o}{=}\PYG{+w}{ }\PYG{n}{GapBuffer}\PYG{p}{::}\PYG{n}{new}\PYG{p}{(}\PYG{p}{)}\PYG{p}{;}
\PYG{n}{buf}\PYG{p}{.}\PYG{n}{insert\PYGZus{}iter}\PYG{p}{(}\PYG{l+s}{\PYGZdq{}}\PYG{l+s}{Lord of the Rings}\PYG{l+s}{\PYGZdq{}}\PYG{p}{.}\PYG{n}{chars}\PYG{p}{(}\PYG{p}{)}\PYG{p}{)}\PYG{p}{;}
\PYG{n}{buf}\PYG{p}{.}\PYG{n}{set\PYGZus{}position}\PYG{p}{(}\PYG{l+m+mi}{12}\PYG{p}{)}\PYG{p}{;}
\end{Verbatim}

运行上述代码后，缓冲区如图 22-6 所示。

\myFigure{images/ch22/fig-22-6.jpeg}{图 22-6：包含一些文本的间隙缓冲区}

插入就是要用新文本填补间隙。下面这段代码添加了一个单词并破坏了原句要表达的意思：

\begin{Verbatim}[commandchars=\\\{\},breaklines=true,breakanywhere=true,fontsize=\codefont,xleftmargin=17bp]
\PYG{n}{buf}\PYG{p}{.}\PYG{n}{insert\PYGZus{}iter}\PYG{p}{(}\PYG{l+s}{\PYGZdq{}}\PYG{l+s}{Onion }\PYG{l+s}{\PYGZdq{}}\PYG{p}{.}\PYG{n}{chars}\PYG{p}{(}\PYG{p}{)}\PYG{p}{)}\PYG{p}{;}
\end{Verbatim}

这会导致如图 22-7 所示的状态。

\myFigure{images/ch22/fig-22-7.jpeg}{图 22-7：包含更多文本的间隙缓冲区}

下面是我们的 \texttt{GapBuffer} 类型：

\begin{Verbatim}[commandchars=\\\{\},breaklines=true,breakanywhere=true,fontsize=\codefont,xleftmargin=17bp]
\PYG{k}{use}\PYG{+w}{ }\PYG{n}{std}\PYG{p}{;}
\PYG{k}{use}\PYG{+w}{ }\PYG{n}{std}\PYG{p}{::}\PYG{n}{ops}\PYG{p}{::}\PYG{n}{Range}\PYG{p}{;}

\PYG{k}{pub}\PYG{+w}{ }\PYG{k}{struct}\PYG{+w}{ }\PYG{n+nc}{GapBuffer}\PYG{o}{\PYGZlt{}}\PYG{n}{T}\PYG{o}{\PYGZgt{}}\PYG{+w}{ }\PYG{p}{\PYGZob{}}
\PYG{+w}{    }\PYG{c+c1}{// 元素的存储区。这个存储区具有我们需要的容量，但它的长度始终保持为0。}
\PYG{+w}{    }\PYG{c+c1}{// GapBuffer会将其元素和间隙放入此`Vec`的“未使用”容量中}
\PYG{+w}{    }\PYG{n}{storage}\PYG{p}{:}\PYG{+w}{ }\PYG{n+nb}{Vec}\PYG{o}{\PYGZlt{}}\PYG{n}{T}\PYG{o}{\PYGZgt{}}\PYG{p}{,}

\PYG{+w}{    }\PYG{c+c1}{// `storage`中间未初始化元素的范围}
\PYG{+w}{    }\PYG{c+c1}{// 这个范围前后的元素始终是已初始化的}
\PYG{+w}{    }\PYG{n}{gap}\PYG{p}{:}\PYG{+w}{ }\PYG{n+nc}{Range}\PYG{o}{\PYGZlt{}}\PYG{k+kt}{usize}\PYG{o}{\PYGZgt{}}
\PYG{p}{\PYGZcb{}}
\end{Verbatim}

\texttt{GapBuffer} 会以一种奇怪的方式使用它的 \texttt{storage} 字段。\footnote{使用编译器内部 \texttt{alloc} crate 中的 \texttt{RawVec} 类型能更好地处理这个问题，但那个 crate 还不稳定。}它实际上从未在向量中存储任何元素（不过这么说也不太准确），而只是简单地调用 \texttt{Vec::with\_capacity(n)} 来获取一块足够大的内存以容纳 \texttt{n} 值，通过向量的 \texttt{as\_ptr} 方法和 \texttt{as\_mut\_ptr} 方法获得指向该内存的裸指针，然后直接将该缓冲区用于自己的目的。向量的长度始终保持为 0。当 \texttt{Vec} 被丢弃时，\texttt{Vec} 不会尝试释放自己的元素（因为它不认为自己有任何元素），而只会释放内存块。这正是 \texttt{GapBuffer} 想要的行为，它有自己的 \texttt{Drop} 实现，知道有效元素在哪里并能正确地丢弃它们。

\texttt{GapBuffer} 中最简单的方法正如你所预期的：

\begin{Verbatim}[commandchars=\\\{\},breaklines=true,breakanywhere=true,fontsize=\codefont,xleftmargin=17bp]
\PYG{k}{impl}\PYG{o}{\PYGZlt{}}\PYG{n}{T}\PYG{o}{\PYGZgt{}}\PYG{+w}{ }\PYG{n}{GapBuffer}\PYG{o}{\PYGZlt{}}\PYG{n}{T}\PYG{o}{\PYGZgt{}}\PYG{+w}{ }\PYG{p}{\PYGZob{}}
\PYG{+w}{    }\PYG{k}{pub}\PYG{+w}{ }\PYG{k}{fn}\PYG{+w}{ }\PYG{n+nf}{new}\PYG{p}{(}\PYG{p}{)}\PYG{+w}{ }\PYG{p}{\PYGZhy{}\PYGZgt{}}\PYG{+w}{ }\PYG{n+nc}{GapBuffer}\PYG{o}{\PYGZlt{}}\PYG{n}{T}\PYG{o}{\PYGZgt{}}\PYG{+w}{ }\PYG{p}{\PYGZob{}}
\PYG{+w}{        }\PYG{n}{GapBuffer}\PYG{+w}{ }\PYG{p}{\PYGZob{}}\PYG{+w}{ }\PYG{n}{storage}\PYG{p}{:}\PYG{+w}{ }\PYG{n+nb}{Vec}\PYG{p}{::}\PYG{n}{new}\PYG{p}{(}\PYG{p}{)}\PYG{p}{,}\PYG{+w}{ }\PYG{n}{gap}\PYG{p}{:}\PYG{+w}{ }\PYG{l+m+mi}{0}\PYG{o}{..}\PYG{l+m+mi}{0}\PYG{+w}{ }\PYG{p}{\PYGZcb{}}
\PYG{+w}{    }\PYG{p}{\PYGZcb{}}

\PYG{+w}{    }\PYG{l+s+sd}{/// 返回在不重新分配的情况下这个GapBuffer可以容纳的元素数}
\PYG{+w}{    }\PYG{k}{pub}\PYG{+w}{ }\PYG{k}{fn}\PYG{+w}{ }\PYG{n+nf}{capacity}\PYG{p}{(}\PYG{o}{\PYGZam{}}\PYG{n+nb+bp}{self}\PYG{p}{)}\PYG{+w}{ }\PYG{p}{\PYGZhy{}\PYGZgt{}}\PYG{+w}{ }\PYG{k+kt}{usize}\PYG{+w}{ }\PYG{p}{\PYGZob{}}
\PYG{+w}{        }\PYG{n+nb+bp}{self}\PYG{p}{.}\PYG{n}{storage}\PYG{p}{.}\PYG{n}{capacity}\PYG{p}{(}\PYG{p}{)}
\PYG{+w}{    }\PYG{p}{\PYGZcb{}}

\PYG{+w}{    }\PYG{l+s+sd}{/// 返回这个GapBuffer当前包含的元素数}
\PYG{+w}{    }\PYG{k}{pub}\PYG{+w}{ }\PYG{k}{fn}\PYG{+w}{ }\PYG{n+nf}{len}\PYG{p}{(}\PYG{o}{\PYGZam{}}\PYG{n+nb+bp}{self}\PYG{p}{)}\PYG{+w}{ }\PYG{p}{\PYGZhy{}\PYGZgt{}}\PYG{+w}{ }\PYG{k+kt}{usize}\PYG{+w}{ }\PYG{p}{\PYGZob{}}
\PYG{+w}{        }\PYG{n+nb+bp}{self}\PYG{p}{.}\PYG{n}{capacity}\PYG{p}{(}\PYG{p}{)}\PYG{+w}{ }\PYG{o}{\PYGZhy{}}\PYG{+w}{ }\PYG{n+nb+bp}{self}\PYG{p}{.}\PYG{n}{gap}\PYG{p}{.}\PYG{n}{len}\PYG{p}{(}\PYG{p}{)}
\PYG{+w}{    }\PYG{p}{\PYGZcb{}}

\PYG{+w}{    }\PYG{l+s+sd}{/// 返回当前插入点}
\PYG{+w}{    }\PYG{k}{pub}\PYG{+w}{ }\PYG{k}{fn}\PYG{+w}{ }\PYG{n+nf}{position}\PYG{p}{(}\PYG{o}{\PYGZam{}}\PYG{n+nb+bp}{self}\PYG{p}{)}\PYG{+w}{ }\PYG{p}{\PYGZhy{}\PYGZgt{}}\PYG{+w}{ }\PYG{k+kt}{usize}\PYG{+w}{ }\PYG{p}{\PYGZob{}}
\PYG{+w}{        }\PYG{n+nb+bp}{self}\PYG{p}{.}\PYG{n}{gap}\PYG{p}{.}\PYG{n}{start}
\PYG{+w}{    }\PYG{p}{\PYGZcb{}}

\PYG{+w}{    }\PYG{o}{..}\PYG{p}{.}
\PYG{p}{\PYGZcb{}}
\end{Verbatim}

它为后面的很多函数提供了一个工具方法，简化了那些函数的实现。该工具方法会返回指向缓冲区中给定索引处元素的裸指针。为了满足 Rust 的要求，需要为 \texttt{mut} 指针和 \texttt{const} 指针分别定义一个方法。与前面的方法不同，这些方法都不是公共的。继续看这个 \texttt{impl} 块：

\begin{Verbatim}[commandchars=\\\{\},breaklines=true,breakanywhere=true,fontsize=\codefont,xleftmargin=17bp]
\PYG{l+s+sd}{/// 返回底层存储中第`index`个元素的指针，不考虑间隙}
\PYG{l+s+sd}{///}
\PYG{l+s+sd}{/// 安全性: `index`必须是`self.storage`中的有效索引}
\PYG{k}{unsafe}\PYG{+w}{ }\PYG{k}{fn}\PYG{+w}{ }\PYG{n+nf}{space}\PYG{p}{(}\PYG{o}{\PYGZam{}}\PYG{n+nb+bp}{self}\PYG{p}{,}\PYG{+w}{ }\PYG{n}{index}\PYG{p}{:}\PYG{+w}{ }\PYG{k+kt}{usize}\PYG{p}{)}\PYG{+w}{ }\PYG{p}{\PYGZhy{}\PYGZgt{}}\PYG{+w}{ }\PYG{o}{*}\PYG{k}{const}\PYG{+w}{ }\PYG{n}{T}\PYG{+w}{ }\PYG{p}{\PYGZob{}}
\PYG{+w}{    }\PYG{n+nb+bp}{self}\PYG{p}{.}\PYG{n}{storage}\PYG{p}{.}\PYG{n}{as\PYGZus{}ptr}\PYG{p}{(}\PYG{p}{)}\PYG{p}{.}\PYG{n}{offset}\PYG{p}{(}\PYG{n}{index}\PYG{+w}{ }\PYG{k}{as}\PYG{+w}{ }\PYG{k+kt}{isize}\PYG{p}{)}
\PYG{p}{\PYGZcb{}}

\PYG{l+s+sd}{/// 返回底层存储中第`index`个元素的可变指针，不考虑间隙}
\PYG{l+s+sd}{///}
\PYG{l+s+sd}{/// 安全性：`index`必须是`self.storage`中的有效索引}
\PYG{k}{unsafe}\PYG{+w}{ }\PYG{k}{fn}\PYG{+w}{ }\PYG{n+nf}{space\PYGZus{}mut}\PYG{p}{(}\PYG{o}{\PYGZam{}}\PYG{k}{mut}\PYG{+w}{ }\PYG{n+nb+bp}{self}\PYG{p}{,}\PYG{+w}{ }\PYG{n}{index}\PYG{p}{:}\PYG{+w}{ }\PYG{k+kt}{usize}\PYG{p}{)}\PYG{+w}{ }\PYG{p}{\PYGZhy{}\PYGZgt{}}\PYG{+w}{ }\PYG{o}{*}\PYG{k}{mut}\PYG{+w}{ }\PYG{n}{T}\PYG{+w}{ }\PYG{p}{\PYGZob{}}
\PYG{+w}{    }\PYG{n+nb+bp}{self}\PYG{p}{.}\PYG{n}{storage}\PYG{p}{.}\PYG{n}{as\PYGZus{}mut\PYGZus{}ptr}\PYG{p}{(}\PYG{p}{)}\PYG{p}{.}\PYG{n}{offset}\PYG{p}{(}\PYG{n}{index}\PYG{+w}{ }\PYG{k}{as}\PYG{+w}{ }\PYG{k+kt}{isize}\PYG{p}{)}
\PYG{p}{\PYGZcb{}}
\end{Verbatim}

要找到给定索引处的元素，就必须考虑该索引是落在间隙之前还是之后，并适当调整：

\begin{Verbatim}[commandchars=\\\{\},breaklines=true,breakanywhere=true,fontsize=\codefont,xleftmargin=17bp]
\PYG{l+s+sd}{/// 返回缓冲区中第`index`个元素的偏移量，并将间隙考虑在内。}
\PYG{l+s+sd}{/// 这个方法不检查索引是否在范围内，但永远不会返回间隙中的索引}
\PYG{k}{fn}\PYG{+w}{ }\PYG{n+nf}{index\PYGZus{}to\PYGZus{}raw}\PYG{p}{(}\PYG{o}{\PYGZam{}}\PYG{n+nb+bp}{self}\PYG{p}{,}\PYG{+w}{ }\PYG{n}{index}\PYG{p}{:}\PYG{+w}{ }\PYG{k+kt}{usize}\PYG{p}{)}\PYG{+w}{ }\PYG{p}{\PYGZhy{}\PYGZgt{}}\PYG{+w}{ }\PYG{k+kt}{usize}\PYG{+w}{ }\PYG{p}{\PYGZob{}}
\PYG{+w}{    }\PYG{k}{if}\PYG{+w}{ }\PYG{n}{index}\PYG{+w}{ }\PYG{o}{\PYGZlt{}}\PYG{+w}{ }\PYG{n+nb+bp}{self}\PYG{p}{.}\PYG{n}{gap}\PYG{p}{.}\PYG{n}{start}\PYG{+w}{ }\PYG{p}{\PYGZob{}}
\PYG{+w}{        }\PYG{n}{index}
\PYG{+w}{    }\PYG{p}{\PYGZcb{}}\PYG{+w}{ }\PYG{k}{else}\PYG{+w}{ }\PYG{p}{\PYGZob{}}
\PYG{+w}{        }\PYG{n}{index}\PYG{+w}{ }\PYG{o}{+}\PYG{+w}{ }\PYG{n+nb+bp}{self}\PYG{p}{.}\PYG{n}{gap}\PYG{p}{.}\PYG{n}{len}\PYG{p}{(}\PYG{p}{)}
\PYG{+w}{    }\PYG{p}{\PYGZcb{}}
\PYG{p}{\PYGZcb{}}

\PYG{l+s+sd}{/// 返回对第`index`个元素的引用，如果`index`超出了范围，则返回`None`}
\PYG{k}{pub}\PYG{+w}{ }\PYG{k}{fn}\PYG{+w}{ }\PYG{n+nf}{get}\PYG{p}{(}\PYG{o}{\PYGZam{}}\PYG{n+nb+bp}{self}\PYG{p}{,}\PYG{+w}{ }\PYG{n}{index}\PYG{p}{:}\PYG{+w}{ }\PYG{k+kt}{usize}\PYG{p}{)}\PYG{+w}{ }\PYG{p}{\PYGZhy{}\PYGZgt{}}\PYG{+w}{ }\PYG{n+nb}{Option}\PYG{o}{\PYGZlt{}}\PYG{o}{\PYGZam{}}\PYG{n}{T}\PYG{o}{\PYGZgt{}}\PYG{+w}{ }\PYG{p}{\PYGZob{}}
\PYG{+w}{    }\PYG{k+kd}{let}\PYG{+w}{ }\PYG{n}{raw}\PYG{+w}{ }\PYG{o}{=}\PYG{+w}{ }\PYG{n+nb+bp}{self}\PYG{p}{.}\PYG{n}{index\PYGZus{}to\PYGZus{}raw}\PYG{p}{(}\PYG{n}{index}\PYG{p}{)}\PYG{p}{;}
\PYG{+w}{    }\PYG{k}{if}\PYG{+w}{ }\PYG{n}{raw}\PYG{+w}{ }\PYG{o}{\PYGZlt{}}\PYG{+w}{ }\PYG{n+nb+bp}{self}\PYG{p}{.}\PYG{n}{capacity}\PYG{p}{(}\PYG{p}{)}\PYG{+w}{ }\PYG{p}{\PYGZob{}}
\PYG{+w}{        }\PYG{k}{unsafe}\PYG{+w}{ }\PYG{p}{\PYGZob{}}
\PYG{+w}{            }\PYG{c+c1}{// 刚刚针对self.capacity()检查过`raw`，而index\PYGZus{}to\PYGZus{}raw}
\PYG{+w}{            }\PYG{c+c1}{// 跳过了间隙，所以这是安全的}
\PYG{+w}{            }\PYG{n+nb}{Some}\PYG{p}{(}\PYG{o}{\PYGZam{}}\PYG{o}{*}\PYG{n+nb+bp}{self}\PYG{p}{.}\PYG{n}{space}\PYG{p}{(}\PYG{n}{raw}\PYG{p}{)}\PYG{p}{)}
\PYG{+w}{        }\PYG{p}{\PYGZcb{}}
\PYG{+w}{    }\PYG{p}{\PYGZcb{}}\PYG{+w}{ }\PYG{k}{else}\PYG{+w}{ }\PYG{p}{\PYGZob{}}
\PYG{+w}{        }\PYG{n+nb}{None}
\PYG{+w}{    }\PYG{p}{\PYGZcb{}}
\PYG{p}{\PYGZcb{}}
\end{Verbatim}

当开始在缓冲区的不同部分进行插入和删除时，需要将间隙移动到新位置。向右移动间隙就要向左移动元素，反之亦然，这就像水平仪中的气泡向一个方向移动而液体会向另一个方向移动一样：

\begin{Verbatim}[commandchars=\\\{\},breaklines=true,breakanywhere=true,fontsize=\codefont,xleftmargin=17bp]
\PYG{l+s+sd}{/// 将当前插入点设置为`pos`。如果`pos`越界，就panic}
\PYG{k}{pub}\PYG{+w}{ }\PYG{k}{fn}\PYG{+w}{ }\PYG{n+nf}{set\PYGZus{}position}\PYG{p}{(}\PYG{o}{\PYGZam{}}\PYG{k}{mut}\PYG{+w}{ }\PYG{n+nb+bp}{self}\PYG{p}{,}\PYG{+w}{ }\PYG{n}{pos}\PYG{p}{:}\PYG{+w}{ }\PYG{k+kt}{usize}\PYG{p}{)}\PYG{+w}{ }\PYG{p}{\PYGZob{}}
\PYG{+w}{    }\PYG{k}{if}\PYG{+w}{ }\PYG{n}{pos}\PYG{+w}{ }\PYG{o}{\PYGZgt{}}\PYG{+w}{ }\PYG{n+nb+bp}{self}\PYG{p}{.}\PYG{n}{len}\PYG{p}{(}\PYG{p}{)}\PYG{+w}{ }\PYG{p}{\PYGZob{}}
\PYG{+w}{        }\PYG{n+nf+fm}{panic!}\PYG{p}{(}\PYG{l+s}{\PYGZdq{}}\PYG{l+s}{index \PYGZob{}\PYGZcb{} out of range for GapBuffer}\PYG{l+s}{\PYGZdq{}}\PYG{p}{,}\PYG{+w}{ }\PYG{n}{pos}\PYG{p}{)}\PYG{p}{;}
\PYG{+w}{    }\PYG{p}{\PYGZcb{}}

\PYG{+w}{    }\PYG{k}{unsafe}\PYG{+w}{ }\PYG{p}{\PYGZob{}}
\PYG{+w}{        }\PYG{k+kd}{let}\PYG{+w}{ }\PYG{n}{gap}\PYG{+w}{ }\PYG{o}{=}\PYG{+w}{ }\PYG{n+nb+bp}{self}\PYG{p}{.}\PYG{n}{gap}\PYG{p}{.}\PYG{n}{clone}\PYG{p}{(}\PYG{p}{)}\PYG{p}{;}
\PYG{+w}{        }\PYG{k}{if}\PYG{+w}{ }\PYG{n}{pos}\PYG{+w}{ }\PYG{o}{\PYGZgt{}}\PYG{+w}{ }\PYG{n}{gap}\PYG{p}{.}\PYG{n}{start}\PYG{+w}{ }\PYG{p}{\PYGZob{}}
\PYG{+w}{            }\PYG{c+c1}{// `pos`位于间隙之后。通过将间隙之后的元素移动到间隙之前来向右移动间隙}
\PYG{+w}{            }\PYG{k+kd}{let}\PYG{+w}{ }\PYG{n}{distance}\PYG{+w}{ }\PYG{o}{=}\PYG{+w}{ }\PYG{n}{pos}\PYG{+w}{ }\PYG{o}{\PYGZhy{}}\PYG{+w}{ }\PYG{n}{gap}\PYG{p}{.}\PYG{n}{start}\PYG{p}{;}
\PYG{+w}{            }\PYG{n}{std}\PYG{p}{::}\PYG{n}{ptr}\PYG{p}{::}\PYG{n}{copy}\PYG{p}{(}\PYG{n+nb+bp}{self}\PYG{p}{.}\PYG{n}{space}\PYG{p}{(}\PYG{n}{gap}\PYG{p}{.}\PYG{n}{end}\PYG{p}{)}\PYG{p}{,}
\PYG{+w}{                           }\PYG{n+nb+bp}{self}\PYG{p}{.}\PYG{n}{space\PYGZus{}mut}\PYG{p}{(}\PYG{n}{gap}\PYG{p}{.}\PYG{n}{start}\PYG{p}{)}\PYG{p}{,}
\PYG{+w}{                           }\PYG{n}{distance}\PYG{p}{)}\PYG{p}{;}
\PYG{+w}{        }\PYG{p}{\PYGZcb{}}\PYG{+w}{ }\PYG{k}{else}\PYG{+w}{ }\PYG{k}{if}\PYG{+w}{ }\PYG{n}{pos}\PYG{+w}{ }\PYG{o}{\PYGZlt{}}\PYG{+w}{ }\PYG{n}{gap}\PYG{p}{.}\PYG{n}{start}\PYG{+w}{ }\PYG{p}{\PYGZob{}}
\PYG{+w}{            }\PYG{c+c1}{// `pos`位于间隙之前。通过将间隙之前的元素移动到间隙之后来向左移动间隙}
\PYG{+w}{            }\PYG{k+kd}{let}\PYG{+w}{ }\PYG{n}{distance}\PYG{+w}{ }\PYG{o}{=}\PYG{+w}{ }\PYG{n}{gap}\PYG{p}{.}\PYG{n}{start}\PYG{+w}{ }\PYG{o}{\PYGZhy{}}\PYG{+w}{ }\PYG{n}{pos}\PYG{p}{;}
\PYG{+w}{            }\PYG{n}{std}\PYG{p}{::}\PYG{n}{ptr}\PYG{p}{::}\PYG{n}{copy}\PYG{p}{(}\PYG{n+nb+bp}{self}\PYG{p}{.}\PYG{n}{space}\PYG{p}{(}\PYG{n}{pos}\PYG{p}{)}\PYG{p}{,}
\PYG{+w}{                           }\PYG{n+nb+bp}{self}\PYG{p}{.}\PYG{n}{space\PYGZus{}mut}\PYG{p}{(}\PYG{n}{gap}\PYG{p}{.}\PYG{n}{end}\PYG{+w}{ }\PYG{o}{\PYGZhy{}}\PYG{+w}{ }\PYG{n}{distance}\PYG{p}{)}\PYG{p}{,}
\PYG{+w}{                           }\PYG{n}{distance}\PYG{p}{)}\PYG{p}{;}
\PYG{+w}{        }\PYG{p}{\PYGZcb{}}

\PYG{+w}{        }\PYG{n+nb+bp}{self}\PYG{p}{.}\PYG{n}{gap}\PYG{+w}{ }\PYG{o}{=}\PYG{+w}{ }\PYG{n}{pos}\PYG{+w}{ }\PYG{o}{..}\PYG{+w}{ }\PYG{n}{pos}\PYG{+w}{ }\PYG{o}{+}\PYG{+w}{ }\PYG{n}{gap}\PYG{p}{.}\PYG{n}{len}\PYG{p}{(}\PYG{p}{)}\PYG{p}{;}
\PYG{+w}{    }\PYG{p}{\PYGZcb{}}
\PYG{p}{\PYGZcb{}}
\end{Verbatim}

这个函数使用 \texttt{std::ptr::copy} 方法来平移元素，\texttt{copy} 要求目标是未初始化的并且会让源保持未初始化。源和目标范围可以重叠，\texttt{copy} 会正确处理这种情况。由于间隙是调用前尚未初始化的内存，而这个函数会调整间隙的位置以覆盖 \texttt{copy} 腾出的空间，因此可以满足 \texttt{copy} 函数的契约。

元素的插入和移除都比较简单。插入会从间隙中为新元素占用一个空间，而移除会将值移出并扩大间隙以覆盖此值曾占据的空间：

\begin{Verbatim}[commandchars=\\\{\},breaklines=true,breakanywhere=true,fontsize=\codefont,xleftmargin=17bp]
\PYG{l+s+sd}{/// 在当前插入点插入`elt`，并在插入后把插入点后移}
\PYG{k}{pub}\PYG{+w}{ }\PYG{k}{fn}\PYG{+w}{ }\PYG{n+nf}{insert}\PYG{p}{(}\PYG{o}{\PYGZam{}}\PYG{k}{mut}\PYG{+w}{ }\PYG{n+nb+bp}{self}\PYG{p}{,}\PYG{+w}{ }\PYG{n}{elt}\PYG{p}{:}\PYG{+w}{ }\PYG{n+nc}{T}\PYG{p}{)}\PYG{+w}{ }\PYG{p}{\PYGZob{}}
\PYG{+w}{    }\PYG{k}{if}\PYG{+w}{ }\PYG{n+nb+bp}{self}\PYG{p}{.}\PYG{n}{gap}\PYG{p}{.}\PYG{n}{len}\PYG{p}{(}\PYG{p}{)}\PYG{+w}{ }\PYG{o}{=}\PYG{o}{=}\PYG{+w}{ }\PYG{l+m+mi}{0}\PYG{+w}{ }\PYG{p}{\PYGZob{}}
\PYG{+w}{        }\PYG{n+nb+bp}{self}\PYG{p}{.}\PYG{n}{enlarge\PYGZus{}gap}\PYG{p}{(}\PYG{p}{)}\PYG{p}{;}
\PYG{+w}{    }\PYG{p}{\PYGZcb{}}

\PYG{+w}{    }\PYG{k}{unsafe}\PYG{+w}{ }\PYG{p}{\PYGZob{}}
\PYG{+w}{        }\PYG{k+kd}{let}\PYG{+w}{ }\PYG{n}{index}\PYG{+w}{ }\PYG{o}{=}\PYG{+w}{ }\PYG{n+nb+bp}{self}\PYG{p}{.}\PYG{n}{gap}\PYG{p}{.}\PYG{n}{start}\PYG{p}{;}
\PYG{+w}{        }\PYG{n}{std}\PYG{p}{::}\PYG{n}{ptr}\PYG{p}{::}\PYG{n}{write}\PYG{p}{(}\PYG{n+nb+bp}{self}\PYG{p}{.}\PYG{n}{space\PYGZus{}mut}\PYG{p}{(}\PYG{n}{index}\PYG{p}{)}\PYG{p}{,}\PYG{+w}{ }\PYG{n}{elt}\PYG{p}{)}\PYG{p}{;}
\PYG{+w}{    }\PYG{p}{\PYGZcb{}}
\PYG{+w}{    }\PYG{n+nb+bp}{self}\PYG{p}{.}\PYG{n}{gap}\PYG{p}{.}\PYG{n}{start}\PYG{+w}{ }\PYG{o}{+}\PYG{o}{=}\PYG{+w}{ }\PYG{l+m+mi}{1}\PYG{p}{;}
\PYG{p}{\PYGZcb{}}

\PYG{l+s+sd}{/// 在当前插入点插入`iter`生成的元素，并在插入后把插入点后移}
\PYG{k}{pub}\PYG{+w}{ }\PYG{k}{fn}\PYG{+w}{ }\PYG{n+nf}{insert\PYGZus{}iter}\PYG{o}{\PYGZlt{}}\PYG{n}{I}\PYG{o}{\PYGZgt{}}\PYG{p}{(}\PYG{o}{\PYGZam{}}\PYG{k}{mut}\PYG{+w}{ }\PYG{n+nb+bp}{self}\PYG{p}{,}\PYG{+w}{ }\PYG{n}{iterable}\PYG{p}{:}\PYG{+w}{ }\PYG{n+nc}{I}\PYG{p}{)}
\PYG{+w}{    }\PYG{k}{where}\PYG{+w}{ }\PYG{n}{I}\PYG{p}{:}\PYG{+w}{ }\PYG{n+nb}{IntoIterator}\PYG{o}{\PYGZlt{}}\PYG{n}{Item}\PYG{o}{=}\PYG{n}{T}\PYG{o}{\PYGZgt{}}
\PYG{p}{\PYGZob{}}
\PYG{+w}{    }\PYG{k}{for}\PYG{+w}{ }\PYG{n}{item}\PYG{+w}{ }\PYG{k}{in}\PYG{+w}{ }\PYG{n}{iterable}\PYG{+w}{ }\PYG{p}{\PYGZob{}}
\PYG{+w}{        }\PYG{n+nb+bp}{self}\PYG{p}{.}\PYG{n}{insert}\PYG{p}{(}\PYG{n}{item}\PYG{p}{)}
\PYG{+w}{    }\PYG{p}{\PYGZcb{}}
\PYG{p}{\PYGZcb{}}

\PYG{l+s+sd}{/// 删除插入点之后的元素并返回它，如果插入点位于GapBuffer的末尾，则返回`None`}
\PYG{k}{pub}\PYG{+w}{ }\PYG{k}{fn}\PYG{+w}{ }\PYG{n+nf}{remove}\PYG{p}{(}\PYG{o}{\PYGZam{}}\PYG{k}{mut}\PYG{+w}{ }\PYG{n+nb+bp}{self}\PYG{p}{)}\PYG{+w}{ }\PYG{p}{\PYGZhy{}\PYGZgt{}}\PYG{+w}{ }\PYG{n+nb}{Option}\PYG{o}{\PYGZlt{}}\PYG{n}{T}\PYG{o}{\PYGZgt{}}\PYG{+w}{ }\PYG{p}{\PYGZob{}}
\PYG{+w}{    }\PYG{k}{if}\PYG{+w}{ }\PYG{n+nb+bp}{self}\PYG{p}{.}\PYG{n}{gap}\PYG{p}{.}\PYG{n}{end}\PYG{+w}{ }\PYG{o}{=}\PYG{o}{=}\PYG{+w}{ }\PYG{n+nb+bp}{self}\PYG{p}{.}\PYG{n}{capacity}\PYG{p}{(}\PYG{p}{)}\PYG{+w}{ }\PYG{p}{\PYGZob{}}
\PYG{+w}{        }\PYG{k}{return}\PYG{+w}{ }\PYG{n+nb}{None}\PYG{p}{;}
\PYG{+w}{    }\PYG{p}{\PYGZcb{}}

\PYG{+w}{    }\PYG{k+kd}{let}\PYG{+w}{ }\PYG{n}{element}\PYG{+w}{ }\PYG{o}{=}\PYG{+w}{ }\PYG{k}{unsafe}\PYG{+w}{ }\PYG{p}{\PYGZob{}}
\PYG{+w}{        }\PYG{n}{std}\PYG{p}{::}\PYG{n}{ptr}\PYG{p}{::}\PYG{n}{read}\PYG{p}{(}\PYG{n+nb+bp}{self}\PYG{p}{.}\PYG{n}{space}\PYG{p}{(}\PYG{n+nb+bp}{self}\PYG{p}{.}\PYG{n}{gap}\PYG{p}{.}\PYG{n}{end}\PYG{p}{)}\PYG{p}{)}
\PYG{+w}{    }\PYG{p}{\PYGZcb{}}\PYG{p}{;}
\PYG{+w}{    }\PYG{n+nb+bp}{self}\PYG{p}{.}\PYG{n}{gap}\PYG{p}{.}\PYG{n}{end}\PYG{+w}{ }\PYG{o}{+}\PYG{o}{=}\PYG{+w}{ }\PYG{l+m+mi}{1}\PYG{p}{;}
\PYG{+w}{    }\PYG{n+nb}{Some}\PYG{p}{(}\PYG{n}{element}\PYG{p}{)}
\PYG{p}{\PYGZcb{}}
\end{Verbatim}

与 \texttt{Vec} 使用 \texttt{std::ptr::write} 进行 \texttt{push} 和使用 \texttt{std::ptr::read} 进行 \texttt{pop} 的方式类似，\texttt{GapBuffer} 使用 \texttt{write} 进行 \texttt{insert}，使用 \texttt{read} 进行 \texttt{remove}。与 \texttt{Vec} 必须调整其长度以维持已初始化元素和空闲容量之间的边界一样，\texttt{GapBuffer} 也会调整其间隙。

填补此间隙后，\texttt{insert} 方法必须扩大缓冲区以获得更多可用空间。\texttt{enlarge\_gap} 方法（\texttt{impl} 块中的最后一个）会处理这个问题：

\begin{Verbatim}[commandchars=\\\{\},breaklines=true,breakanywhere=true,fontsize=\codefont,xleftmargin=17bp]
\PYG{l+s+sd}{/// 将`self.storage`的容量翻倍}
\PYG{k}{fn}\PYG{+w}{ }\PYG{n+nf}{enlarge\PYGZus{}gap}\PYG{p}{(}\PYG{o}{\PYGZam{}}\PYG{k}{mut}\PYG{+w}{ }\PYG{n+nb+bp}{self}\PYG{p}{)}\PYG{+w}{ }\PYG{p}{\PYGZob{}}
\PYG{+w}{    }\PYG{k+kd}{let}\PYG{+w}{ }\PYG{k}{mut}\PYG{+w}{ }\PYG{n}{new\PYGZus{}capacity}\PYG{+w}{ }\PYG{o}{=}\PYG{+w}{ }\PYG{n+nb+bp}{self}\PYG{p}{.}\PYG{n}{capacity}\PYG{p}{(}\PYG{p}{)}\PYG{+w}{ }\PYG{o}{*}\PYG{+w}{ }\PYG{l+m+mi}{2}\PYG{p}{;}
\PYG{+w}{    }\PYG{k}{if}\PYG{+w}{ }\PYG{n}{new\PYGZus{}capacity}\PYG{+w}{ }\PYG{o}{=}\PYG{o}{=}\PYG{+w}{ }\PYG{l+m+mi}{0}\PYG{+w}{ }\PYG{p}{\PYGZob{}}
\PYG{+w}{        }\PYG{c+c1}{// 现有向量是空的}
\PYG{+w}{        }\PYG{c+c1}{// 选择一个合理的初始容量}
\PYG{+w}{        }\PYG{n}{new\PYGZus{}capacity}\PYG{+w}{ }\PYG{o}{=}\PYG{+w}{ }\PYG{l+m+mi}{4}\PYG{p}{;}
\PYG{+w}{    }\PYG{p}{\PYGZcb{}}

\PYG{+w}{    }\PYG{c+c1}{// 我们不知道调整Vec的大小会对其“（表面看）未使用的”容量}
\PYG{+w}{    }\PYG{c+c1}{// 有何影响，所以只好创建一个新向量并把元素移了过去}
\PYG{+w}{    }\PYG{k+kd}{let}\PYG{+w}{ }\PYG{k}{mut}\PYG{+w}{ }\PYG{n}{new}\PYG{+w}{ }\PYG{o}{=}\PYG{+w}{ }\PYG{n+nb}{Vec}\PYG{p}{::}\PYG{n}{with\PYGZus{}capacity}\PYG{p}{(}\PYG{n}{new\PYGZus{}capacity}\PYG{p}{)}\PYG{p}{;}
\PYG{+w}{    }\PYG{k+kd}{let}\PYG{+w}{ }\PYG{n}{after\PYGZus{}gap}\PYG{+w}{ }\PYG{o}{=}\PYG{+w}{ }\PYG{n+nb+bp}{self}\PYG{p}{.}\PYG{n}{capacity}\PYG{p}{(}\PYG{p}{)}\PYG{+w}{ }\PYG{o}{\PYGZhy{}}\PYG{+w}{ }\PYG{n+nb+bp}{self}\PYG{p}{.}\PYG{n}{gap}\PYG{p}{.}\PYG{n}{end}\PYG{p}{;}
\PYG{+w}{    }\PYG{k+kd}{let}\PYG{+w}{ }\PYG{n}{new\PYGZus{}gap}\PYG{+w}{ }\PYG{o}{=}\PYG{+w}{ }\PYG{n+nb+bp}{self}\PYG{p}{.}\PYG{n}{gap}\PYG{p}{.}\PYG{n}{start}\PYG{+w}{ }\PYG{o}{..}\PYG{+w}{ }\PYG{n}{new}\PYG{p}{.}\PYG{n}{capacity}\PYG{p}{(}\PYG{p}{)}\PYG{+w}{ }\PYG{o}{\PYGZhy{}}\PYG{+w}{ }\PYG{n}{after\PYGZus{}gap}\PYG{p}{;}

\PYG{+w}{    }\PYG{k}{unsafe}\PYG{+w}{ }\PYG{p}{\PYGZob{}}
\PYG{+w}{        }\PYG{c+c1}{// 移动位于此间隙之前的元素}
\PYG{+w}{        }\PYG{n}{std}\PYG{p}{::}\PYG{n}{ptr}\PYG{p}{::}\PYG{n}{copy\PYGZus{}nonoverlapping}\PYG{p}{(}\PYG{n+nb+bp}{self}\PYG{p}{.}\PYG{n}{space}\PYG{p}{(}\PYG{l+m+mi}{0}\PYG{p}{)}\PYG{p}{,}
\PYG{+w}{                                      }\PYG{n}{new}\PYG{p}{.}\PYG{n}{as\PYGZus{}mut\PYGZus{}ptr}\PYG{p}{(}\PYG{p}{)}\PYG{p}{,}
\PYG{+w}{                                      }\PYG{n+nb+bp}{self}\PYG{p}{.}\PYG{n}{gap}\PYG{p}{.}\PYG{n}{start}\PYG{p}{)}\PYG{p}{;}

\PYG{+w}{        }\PYG{c+c1}{// 移动位于此间隙之后的元素}
\PYG{+w}{        }\PYG{k+kd}{let}\PYG{+w}{ }\PYG{n}{new\PYGZus{}gap\PYGZus{}end}\PYG{+w}{ }\PYG{o}{=}\PYG{+w}{ }\PYG{n}{new}\PYG{p}{.}\PYG{n}{as\PYGZus{}mut\PYGZus{}ptr}\PYG{p}{(}\PYG{p}{)}\PYG{p}{.}\PYG{n}{offset}\PYG{p}{(}\PYG{n}{new\PYGZus{}gap}\PYG{p}{.}\PYG{n}{end}\PYG{+w}{ }\PYG{k}{as}\PYG{+w}{ }\PYG{k+kt}{isize}\PYG{p}{)}\PYG{p}{;}
\PYG{+w}{        }\PYG{n}{std}\PYG{p}{::}\PYG{n}{ptr}\PYG{p}{::}\PYG{n}{copy\PYGZus{}nonoverlapping}\PYG{p}{(}\PYG{n+nb+bp}{self}\PYG{p}{.}\PYG{n}{space}\PYG{p}{(}\PYG{n+nb+bp}{self}\PYG{p}{.}\PYG{n}{gap}\PYG{p}{.}\PYG{n}{end}\PYG{p}{)}\PYG{p}{,}
\PYG{+w}{                                      }\PYG{n}{new\PYGZus{}gap\PYGZus{}end}\PYG{p}{,}
\PYG{+w}{                                      }\PYG{n}{after\PYGZus{}gap}\PYG{p}{)}\PYG{p}{;}
\PYG{+w}{    }\PYG{p}{\PYGZcb{}}

\PYG{+w}{    }\PYG{c+c1}{// 这会释放旧的Vec，但不会丢弃任何元素，因为此Vec的长度为0}
\PYG{+w}{    }\PYG{n+nb+bp}{self}\PYG{p}{.}\PYG{n}{storage}\PYG{+w}{ }\PYG{o}{=}\PYG{+w}{ }\PYG{n}{new}\PYG{p}{;}
\PYG{+w}{    }\PYG{n+nb+bp}{self}\PYG{p}{.}\PYG{n}{gap}\PYG{+w}{ }\PYG{o}{=}\PYG{+w}{ }\PYG{n}{new\PYGZus{}gap}\PYG{p}{;}
\PYG{p}{\PYGZcb{}}
\end{Verbatim}

\texttt{set\_position} 必须使用 \texttt{copy} 在间隙中来回移动元素，\texttt{enlarge\_gap} 则可以使用 \texttt{copy\_nonoverlapping}，因为它会将元素移动到一个全新的缓冲区。

将新向量转移给 \texttt{self.storage} 会丢弃旧向量。由于旧向量的长度为 0，它认为自己没有要丢弃的元素，因此只释放了自己的缓冲区。巧妙的是，\texttt{copy\_nonoverlapping} 也有把源变成未初始化状态的语义，因此旧向量的做法恰巧是正确的：现在所有元素都归新向量所有了。

最后，需要确保丢弃 \texttt{GapBuffer} 也会丢弃它的所有元素：

\begin{Verbatim}[commandchars=\\\{\},breaklines=true,breakanywhere=true,fontsize=\codefont,xleftmargin=17bp]
\PYG{k}{impl}\PYG{o}{\PYGZlt{}}\PYG{n}{T}\PYG{o}{\PYGZgt{}}\PYG{+w}{ }\PYG{n+nb}{Drop}\PYG{+w}{ }\PYG{k}{for}\PYG{+w}{ }\PYG{n}{GapBuffer}\PYG{o}{\PYGZlt{}}\PYG{n}{T}\PYG{o}{\PYGZgt{}}\PYG{+w}{ }\PYG{p}{\PYGZob{}}
\PYG{+w}{    }\PYG{k}{fn}\PYG{+w}{ }\PYG{n+nf}{drop}\PYG{p}{(}\PYG{o}{\PYGZam{}}\PYG{k}{mut}\PYG{+w}{ }\PYG{n+nb+bp}{self}\PYG{p}{)}\PYG{+w}{ }\PYG{p}{\PYGZob{}}
\PYG{+w}{        }\PYG{k}{unsafe}\PYG{+w}{ }\PYG{p}{\PYGZob{}}
\PYG{+w}{            }\PYG{k}{for}\PYG{+w}{ }\PYG{n}{i}\PYG{+w}{ }\PYG{k}{in}\PYG{+w}{ }\PYG{l+m+mi}{0}\PYG{+w}{ }\PYG{o}{..}\PYG{+w}{ }\PYG{n+nb+bp}{self}\PYG{p}{.}\PYG{n}{gap}\PYG{p}{.}\PYG{n}{start}\PYG{+w}{ }\PYG{p}{\PYGZob{}}
\PYG{+w}{                }\PYG{n}{std}\PYG{p}{::}\PYG{n}{ptr}\PYG{p}{::}\PYG{n}{drop\PYGZus{}in\PYGZus{}place}\PYG{p}{(}\PYG{n+nb+bp}{self}\PYG{p}{.}\PYG{n}{space\PYGZus{}mut}\PYG{p}{(}\PYG{n}{i}\PYG{p}{)}\PYG{p}{)}\PYG{p}{;}
\PYG{+w}{            }\PYG{p}{\PYGZcb{}}
\PYG{+w}{            }\PYG{k}{for}\PYG{+w}{ }\PYG{n}{i}\PYG{+w}{ }\PYG{k}{in}\PYG{+w}{ }\PYG{n+nb+bp}{self}\PYG{p}{.}\PYG{n}{gap}\PYG{p}{.}\PYG{n}{end}\PYG{+w}{ }\PYG{o}{..}\PYG{+w}{ }\PYG{n+nb+bp}{self}\PYG{p}{.}\PYG{n}{capacity}\PYG{p}{(}\PYG{p}{)}\PYG{+w}{ }\PYG{p}{\PYGZob{}}
\PYG{+w}{                }\PYG{n}{std}\PYG{p}{::}\PYG{n}{ptr}\PYG{p}{::}\PYG{n}{drop\PYGZus{}in\PYGZus{}place}\PYG{p}{(}\PYG{n+nb+bp}{self}\PYG{p}{.}\PYG{n}{space\PYGZus{}mut}\PYG{p}{(}\PYG{n}{i}\PYG{p}{)}\PYG{p}{)}\PYG{p}{;}
\PYG{+w}{            }\PYG{p}{\PYGZcb{}}
\PYG{+w}{        }\PYG{p}{\PYGZcb{}}
\PYG{+w}{    }\PYG{p}{\PYGZcb{}}
\PYG{p}{\PYGZcb{}}
\end{Verbatim}

这些元素都位于间隙前后，因此需要遍历每个区域并使用 \texttt{std::ptr::drop\_in\_place} 函数丢弃每个元素。\texttt{drop\_in\_place} 函数是一个行为类似于 \texttt{drop(std::ptr::read(ptr))} 的实用程序，但不会“费心”地将值转移给其调用者（因此适用于无固定大小类型）。就像在 \texttt{enlarge\_gap} 中一样，当向量 \texttt{self.storage} 被丢弃时，它的缓冲区实际上是未初始化的。

与本章展示过的其他类型一样，\texttt{GapBuffer} 会确保自己的不变条件足以遵守所使用的每个不安全特性的契约，因此它的所有公共方法都不需要标记为不安全。\texttt{GapBuffer} 为无法用安全代码高效编写的特性实现了一个安全的接口。

\subsection{不安全代码中的 panic 安全性}

在 Rust 中，panic 通常不会导致未定义行为，\texttt{panic!} 宏并不是不安全特性。但是，当你决定使用不安全代码时，就得考虑 panic 安全性的问题了。

考虑 22.8.7 节中的 \texttt{GapBuffer::remove} 方法：

\begin{Verbatim}[commandchars=\\\{\},breaklines=true,breakanywhere=true,fontsize=\codefont,xleftmargin=17bp]
\PYG{k}{pub}\PYG{+w}{ }\PYG{k}{fn}\PYG{+w}{ }\PYG{n+nf}{remove}\PYG{p}{(}\PYG{o}{\PYGZam{}}\PYG{k}{mut}\PYG{+w}{ }\PYG{n+nb+bp}{self}\PYG{p}{)}\PYG{+w}{ }\PYG{p}{\PYGZhy{}\PYGZgt{}}\PYG{+w}{ }\PYG{n+nb}{Option}\PYG{o}{\PYGZlt{}}\PYG{n}{T}\PYG{o}{\PYGZgt{}}\PYG{+w}{ }\PYG{p}{\PYGZob{}}
\PYG{+w}{    }\PYG{k}{if}\PYG{+w}{ }\PYG{n+nb+bp}{self}\PYG{p}{.}\PYG{n}{gap}\PYG{p}{.}\PYG{n}{end}\PYG{+w}{ }\PYG{o}{=}\PYG{o}{=}\PYG{+w}{ }\PYG{n+nb+bp}{self}\PYG{p}{.}\PYG{n}{capacity}\PYG{p}{(}\PYG{p}{)}\PYG{+w}{ }\PYG{p}{\PYGZob{}}
\PYG{+w}{        }\PYG{k}{return}\PYG{+w}{ }\PYG{n+nb}{None}\PYG{p}{;}
\PYG{+w}{    }\PYG{p}{\PYGZcb{}}

\PYG{+w}{    }\PYG{k+kd}{let}\PYG{+w}{ }\PYG{n}{element}\PYG{+w}{ }\PYG{o}{=}\PYG{+w}{ }\PYG{k}{unsafe}\PYG{+w}{ }\PYG{p}{\PYGZob{}}
\PYG{+w}{        }\PYG{n}{std}\PYG{p}{::}\PYG{n}{ptr}\PYG{p}{::}\PYG{n}{read}\PYG{p}{(}\PYG{n+nb+bp}{self}\PYG{p}{.}\PYG{n}{space}\PYG{p}{(}\PYG{n+nb+bp}{self}\PYG{p}{.}\PYG{n}{gap}\PYG{p}{.}\PYG{n}{end}\PYG{p}{)}\PYG{p}{)}
\PYG{+w}{    }\PYG{p}{\PYGZcb{}}\PYG{p}{;}
\PYG{+w}{    }\PYG{n+nb+bp}{self}\PYG{p}{.}\PYG{n}{gap}\PYG{p}{.}\PYG{n}{end}\PYG{+w}{ }\PYG{o}{+}\PYG{o}{=}\PYG{+w}{ }\PYG{l+m+mi}{1}\PYG{p}{;}
\PYG{+w}{    }\PYG{n+nb}{Some}\PYG{p}{(}\PYG{n}{element}\PYG{p}{)}
\PYG{p}{\PYGZcb{}}
\end{Verbatim}

对 \texttt{read} 的调用会将紧随间隙之后的元素移出缓冲区，留下未初始化的空间。此时，\texttt{GapBuffer} 处于不一致状态：我们打破了间隙外的所有元素都必须是初始化的这个不变条件。幸运的是，下一条语句扩大了间隙以覆盖这个空间，因此当我们返回时，不变条件会再次成立。

但是请考虑一下，如果在调用 \texttt{read} 之后、调整 \texttt{self.gap.end} 之前，此代码尝试使用可能引发 panic 的特性（如索引切片），那么会发生什么呢？在这两个操作之间的任何地方突然退出该方法都会使 \texttt{GapBuffer} 在间隙外留下未初始化的元素。下一次调用 \texttt{remove} 可能会尝试再次读取（\texttt{read}）它，即使仅仅丢弃 \texttt{GapBuffer} 也会尝试运行其 \texttt{drop} 方法。这两者都是未定义行为，因为它们会访问未初始化内存。

类型的方法在执行工作时几乎不可避免地会暂时放松类型的不变条件，然后在返回之前让其回到正轨。方法中间出现的 panic 可能会中断清理过程，使类型处于不一致状态。

如果类型只使用安全代码，那么这种不一致可能会使类型行为诡异，但并不会引入未定义行为。不过使用不安全特性的代码通常会依赖其不变条件来满足这些特性的契约。破坏不变条件会导致契约破损，从而导致未定义行为。

使用不安全特性时，必须特别注意识别这些暂时放松了不变条件的敏感代码区域，并确保它们不会执行任何可能引起 panic 的事情。

\section{用联合体重新解释内存}

虽然 Rust 提供了许多有用的抽象，但最终我们编写的软件只是在操纵字节。联合体是 Rust 最强大的特性之一，用于操纵这些字节并选择如何解释它们。例如，任何 32 位（4 字节）的集合都可以解释为整数或浮点数。任何一种解释都是有效的，不过，将一种数据解释为另一种数据可能会导致其失去意义。

下面是一个用来表示可解释为整数或浮点数的字节集合的联合体：

\begin{Verbatim}[commandchars=\\\{\},breaklines=true,breakanywhere=true,fontsize=\codefont,xleftmargin=17bp]
\PYG{k}{union}\PYG{+w}{ }\PYG{n+nc}{FloatOrInt}\PYG{+w}{ }\PYG{p}{\PYGZob{}}
\PYG{+w}{    }\PYG{n}{f}\PYG{p}{:}\PYG{+w}{ }\PYG{k+kt}{f32}\PYG{p}{,}
\PYG{+w}{    }\PYG{n}{i}\PYG{p}{:}\PYG{+w}{ }\PYG{k+kt}{i32}\PYG{p}{,}
\PYG{p}{\PYGZcb{}}
\end{Verbatim}

这是一个包含两个字段（\texttt{f} 和 \texttt{i}）的联合体。这两个字段可以像结构体的字段一样被赋值，但在构造联合体时，只能选择一个字段，这与结构体不同。结构体的字段会引用内存中的不同位置，而联合体的字段会引用相同位序列的不同解释。赋值给不同的字段只是意味着根据适当的类型覆盖这些位中的一部分或全部。在下面的代码中，\texttt{one} 指向的是单个 32 位内存范围，它首先存储一个按简单整数编码的 \texttt{1}，然后存储一个按 IEEE 754 浮点数编码的 \texttt{1.0}。一旦写入了 \texttt{f}，先前写入的 \texttt{FloatOrInt} 值就会被覆盖：

\begin{Verbatim}[commandchars=\\\{\},breaklines=true,breakanywhere=true,fontsize=\codefont,xleftmargin=17bp]
\PYG{k+kd}{let}\PYG{+w}{ }\PYG{k}{mut}\PYG{+w}{ }\PYG{n}{one}\PYG{+w}{ }\PYG{o}{=}\PYG{+w}{ }\PYG{n}{FloatOrInt}\PYG{+w}{ }\PYG{p}{\PYGZob{}}\PYG{+w}{ }\PYG{n}{i}\PYG{p}{:}\PYG{+w}{ }\PYG{l+m+mi}{1}\PYG{+w}{ }\PYG{p}{\PYGZcb{}}\PYG{p}{;}
\PYG{n+nf+fm}{assert\PYGZus{}eq!}\PYG{p}{(}\PYG{k}{unsafe}\PYG{+w}{ }\PYG{p}{\PYGZob{}}\PYG{+w}{ }\PYG{n}{one}\PYG{p}{.}\PYG{n}{i}\PYG{+w}{ }\PYG{p}{\PYGZcb{}}\PYG{p}{,}\PYG{+w}{ }\PYG{l+m+mh}{0x00\PYGZus{}00\PYGZus{}00\PYGZus{}01}\PYG{p}{)}\PYG{p}{;}
\PYG{n}{one}\PYG{p}{.}\PYG{n}{f}\PYG{+w}{ }\PYG{o}{=}\PYG{+w}{ }\PYG{l+m+mf}{1.0}\PYG{p}{;}
\PYG{n+nf+fm}{assert\PYGZus{}eq!}\PYG{p}{(}\PYG{k}{unsafe}\PYG{+w}{ }\PYG{p}{\PYGZob{}}\PYG{+w}{ }\PYG{n}{one}\PYG{p}{.}\PYG{n}{i}\PYG{+w}{ }\PYG{p}{\PYGZcb{}}\PYG{p}{,}\PYG{+w}{ }\PYG{l+m+mh}{0x3F\PYGZus{}80\PYGZus{}00\PYGZus{}00}\PYG{p}{)}\PYG{p}{;}
\end{Verbatim}

出于同样的原因，联合体的大小会由其最大字段决定。例如，下面这个联合体的大小为 64 位，虽然 \texttt{SmallOrLarge::s} 只是一个 \texttt{bool}：

\begin{Verbatim}[commandchars=\\\{\},breaklines=true,breakanywhere=true,fontsize=\codefont,xleftmargin=17bp]
\PYG{k}{union}\PYG{+w}{ }\PYG{n+nc}{SmallOrLarge}\PYG{+w}{ }\PYG{p}{\PYGZob{}}
\PYG{+w}{    }\PYG{n}{s}\PYG{p}{:}\PYG{+w}{ }\PYG{k+kt}{bool}\PYG{p}{,}
\PYG{+w}{    }\PYG{n}{l}\PYG{p}{:}\PYG{+w}{ }\PYG{k+kt}{u64}
\PYG{p}{\PYGZcb{}}
\end{Verbatim}

虽然构建联合体或对它的字段赋值是完全安全的，但读取联合体的任何字段都是不安全的：

\begin{Verbatim}[commandchars=\\\{\},breaklines=true,breakanywhere=true,fontsize=\codefont,xleftmargin=17bp]
\PYG{k+kd}{let}\PYG{+w}{ }\PYG{n}{u}\PYG{+w}{ }\PYG{o}{=}\PYG{+w}{ }\PYG{n}{SmallOrLarge}\PYG{+w}{ }\PYG{p}{\PYGZob{}}\PYG{+w}{ }\PYG{n}{l}\PYG{p}{:}\PYG{+w}{ }\PYG{l+m+mi}{1337}\PYG{+w}{ }\PYG{p}{\PYGZcb{}}\PYG{p}{;}
\PYG{n+nf+fm}{println!}\PYG{p}{(}\PYG{l+s}{\PYGZdq{}}\PYG{l+s}{\PYGZob{}\PYGZcb{}}\PYG{l+s}{\PYGZdq{}}\PYG{p}{,}\PYG{+w}{ }\PYG{k}{unsafe}\PYG{+w}{ }\PYG{p}{\PYGZob{}}\PYG{n}{u}\PYG{p}{.}\PYG{n}{l}\PYG{p}{\PYGZcb{}}\PYG{p}{)}\PYG{p}{;}\PYG{+w}{ }\PYG{c+c1}{// 打印出1337}
\end{Verbatim}

这是因为与枚举不同，联合体没有标签。编译器不会添加额外的位来区分各种变体。在运行期无法判断 \texttt{SmallOrLarge} 是要该解释为 \texttt{u64} 还是 \texttt{bool}，除非程序有一些额外的上下文。

同时，并没有什么内置手段可以保证给定字段的位模式是有效的。例如，写入 \texttt{SmallOrLarge} 值的 \texttt{l} 字段将覆盖其 \texttt{s} 字段，但它创建的这个位模式并无任何用处，甚至可能都不是有效的 \texttt{bool}。因此，虽然写入联合体字段是安全的，但每次读取都需要 \texttt{unsafe} 代码。仅当 \texttt{s} 字段的各个位可以形成有效的 \texttt{bool} 时才允许从 \texttt{u.s} 读取，否则，这就是未定义行为。

只要把这些限制牢记在心，联合体仍然可以成为临时重新解释某些数据的有用方法，尤其是在针对值的表观而非值本身进行计算时。例如，前面提到的 \texttt{FloatOrInt} 类型可以轻松地打印出浮点数的各个位——即便 \texttt{f32} 没有实现过 \texttt{Binary} 格式化程序：

\begin{Verbatim}[commandchars=\\\{\},breaklines=true,breakanywhere=true,fontsize=\codefont,xleftmargin=17bp]
\PYG{k+kd}{let}\PYG{+w}{ }\PYG{n}{float}\PYG{+w}{ }\PYG{o}{=}\PYG{+w}{ }\PYG{n}{FloatOrInt}\PYG{+w}{ }\PYG{p}{\PYGZob{}}\PYG{+w}{ }\PYG{n}{f}\PYG{p}{:}\PYG{+w}{ }\PYG{l+m+mf}{31337.0}\PYG{+w}{ }\PYG{p}{\PYGZcb{}}\PYG{p}{;}
\PYG{c+c1}{// 打印出1000110111101001101001000000000}
\PYG{n+nf+fm}{println!}\PYG{p}{(}\PYG{l+s}{\PYGZdq{}}\PYG{l+s}{\PYGZob{}:b\PYGZcb{}}\PYG{l+s}{\PYGZdq{}}\PYG{p}{,}\PYG{+w}{ }\PYG{k}{unsafe}\PYG{+w}{ }\PYG{p}{\PYGZob{}}\PYG{+w}{ }\PYG{n}{float}\PYG{p}{.}\PYG{n}{i}\PYG{+w}{ }\PYG{p}{\PYGZcb{}}\PYG{p}{)}\PYG{p}{;}
\end{Verbatim}

虽然几乎可以肯定这些简单示例会在任何版本的编译器上如预期般工作，但并不能保证任何字段都从特定位置开始，除非将某个属性添加到 \texttt{union} 定义中，告诉编译器如何在内存中排布数据。添加属性 \texttt{\#[repr(C)]} 可以保证所有字段都从偏移量 \texttt{0} 而不是编译器喜欢的任何位置开始。有了这个保证，这种改写行为就可以用来提取像整数的符号位这样的单独二进制位了：

\begin{Verbatim}[commandchars=\\\{\},breaklines=true,breakanywhere=true,fontsize=\codefont,xleftmargin=17bp]
\PYG{c+cp}{\PYGZsh{}[}\PYG{c+cp}{repr(C)}\PYG{c+cp}{]}
\PYG{k}{union}\PYG{+w}{ }\PYG{n+nc}{SignExtractor}\PYG{+w}{ }\PYG{p}{\PYGZob{}}
\PYG{+w}{    }\PYG{n}{value}\PYG{p}{:}\PYG{+w}{ }\PYG{k+kt}{i64}\PYG{p}{,}
\PYG{+w}{    }\PYG{n}{bytes}\PYG{p}{:}\PYG{+w}{ }\PYG{p}{[}\PYG{k+kt}{u8}\PYG{p}{;}\PYG{+w}{ }\PYG{l+m+mi}{8}\PYG{p}{]}
\PYG{p}{\PYGZcb{}}

\PYG{k}{fn}\PYG{+w}{ }\PYG{n+nf}{sign}\PYG{p}{(}\PYG{n}{int}\PYG{p}{:}\PYG{+w}{ }\PYG{k+kt}{i64}\PYG{p}{)}\PYG{+w}{ }\PYG{p}{\PYGZhy{}\PYGZgt{}}\PYG{+w}{ }\PYG{k+kt}{bool}\PYG{+w}{ }\PYG{p}{\PYGZob{}}
\PYG{+w}{    }\PYG{k+kd}{let}\PYG{+w}{ }\PYG{n}{se}\PYG{+w}{ }\PYG{o}{=}\PYG{+w}{ }\PYG{n}{SignExtractor}\PYG{+w}{ }\PYG{p}{\PYGZob{}}\PYG{+w}{ }\PYG{n}{value}\PYG{p}{:}\PYG{+w}{ }\PYG{n+nc}{int}\PYG{+w}{ }\PYG{p}{\PYGZcb{}}\PYG{p}{;}
\PYG{+w}{    }\PYG{n+nf+fm}{println!}\PYG{p}{(}\PYG{l+s}{\PYGZdq{}}\PYG{l+s}{\PYGZob{}:b\PYGZcb{} (\PYGZob{}:?\PYGZcb{})}\PYG{l+s}{\PYGZdq{}}\PYG{p}{,}\PYG{+w}{ }\PYG{k}{unsafe}\PYG{+w}{ }\PYG{p}{\PYGZob{}}\PYG{+w}{ }\PYG{n}{se}\PYG{p}{.}\PYG{n}{value}\PYG{+w}{ }\PYG{p}{\PYGZcb{}}\PYG{p}{,}\PYG{+w}{ }\PYG{k}{unsafe}\PYG{+w}{ }\PYG{p}{\PYGZob{}}\PYG{+w}{ }\PYG{n}{se}\PYG{p}{.}\PYG{n}{bytes}\PYG{+w}{ }\PYG{p}{\PYGZcb{}}\PYG{p}{)}\PYG{p}{;}
\PYG{+w}{    }\PYG{k}{unsafe}\PYG{+w}{ }\PYG{p}{\PYGZob{}}\PYG{+w}{ }\PYG{n}{se}\PYG{p}{.}\PYG{n}{bytes}\PYG{p}{[}\PYG{l+m+mi}{7}\PYG{p}{]}\PYG{+w}{ }\PYG{o}{\PYGZgt{}}\PYG{o}{=}\PYG{+w}{ }\PYG{l+m+mb}{0b10000000}\PYG{+w}{ }\PYG{p}{\PYGZcb{}}
\PYG{p}{\PYGZcb{}}

\PYG{n+nf+fm}{assert\PYGZus{}eq!}\PYG{p}{(}\PYG{n}{sign}\PYG{p}{(}\PYG{o}{\PYGZhy{}}\PYG{l+m+mi}{1}\PYG{p}{)}\PYG{p}{,}\PYG{+w}{ }\PYG{k+kc}{true}\PYG{p}{)}\PYG{p}{;}
\PYG{n+nf+fm}{assert\PYGZus{}eq!}\PYG{p}{(}\PYG{n}{sign}\PYG{p}{(}\PYG{l+m+mi}{1}\PYG{p}{)}\PYG{p}{,}\PYG{+w}{ }\PYG{k+kc}{false}\PYG{p}{)}\PYG{p}{;}
\PYG{n+nf+fm}{assert\PYGZus{}eq!}\PYG{p}{(}\PYG{n}{sign}\PYG{p}{(}\PYG{k+kt}{i64}\PYG{p}{::}\PYG{n}{MAX}\PYG{p}{)}\PYG{p}{,}\PYG{+w}{ }\PYG{k+kc}{false}\PYG{p}{)}\PYG{p}{;}
\PYG{n+nf+fm}{assert\PYGZus{}eq!}\PYG{p}{(}\PYG{n}{sign}\PYG{p}{(}\PYG{k+kt}{i64}\PYG{p}{::}\PYG{n}{MIN}\PYG{p}{)}\PYG{p}{,}\PYG{+w}{ }\PYG{k+kc}{true}\PYG{p}{)}\PYG{p}{;}
\end{Verbatim}

在这里，符号位是最高有效字节的最高有效位。因为 x86 处理器是小端（低位在前）的，所以这些字节的顺序是相反的，其最高有效字节不是 \texttt{bytes[0]}，而是 \texttt{bytes[7]}。通常，这不是 Rust 代码必须处理的事情，但是因为这段代码要直接与 \texttt{i64} 的内存中表示法打交道，所以这些底层细节就变得很重要了。

因为不知道该如何丢弃其内容，所以联合体的所有字段都必须是可 \texttt{Copy} 的。但是，如果必须在联合体中存储一个 \texttt{String}，那么也有相应的解决方案，详情请参阅 \texttt{std::mem::ManuallyDrop} 的标准库文档。

\section{匹配联合体}

在 Rust 联合体上进行匹配和在结构体上匹配类似，但每个模式必须指定一个字段：

\begin{Verbatim}[commandchars=\\\{\},breaklines=true,breakanywhere=true,fontsize=\codefont,xleftmargin=17bp]
\PYG{k}{unsafe}\PYG{+w}{ }\PYG{p}{\PYGZob{}}
\PYG{+w}{    }\PYG{k}{match}\PYG{+w}{ }\PYG{n}{u}\PYG{+w}{ }\PYG{p}{\PYGZob{}}
\PYG{+w}{        }\PYG{n}{SmallOrLarge}\PYG{+w}{ }\PYG{p}{\PYGZob{}}\PYG{+w}{ }\PYG{n}{s}\PYG{p}{:}\PYG{+w}{ }\PYG{n+nc}{true}\PYG{+w}{ }\PYG{p}{\PYGZcb{}}\PYG{+w}{ }\PYG{o}{=}\PYG{o}{\PYGZgt{}}\PYG{+w}{ }\PYG{p}{\PYGZob{}}\PYG{+w}{ }\PYG{n+nf+fm}{println!}\PYG{p}{(}\PYG{l+s}{\PYGZdq{}}\PYG{l+s}{boolean true}\PYG{l+s}{\PYGZdq{}}\PYG{p}{)}\PYG{p}{;}\PYG{+w}{ }\PYG{p}{\PYGZcb{}}
\PYG{+w}{        }\PYG{n}{SmallOrLarge}\PYG{+w}{ }\PYG{p}{\PYGZob{}}\PYG{+w}{ }\PYG{n}{l}\PYG{p}{:}\PYG{+w}{ }\PYG{l+m+mi}{2}\PYG{+w}{ }\PYG{p}{\PYGZcb{}}\PYG{+w}{ }\PYG{o}{=}\PYG{o}{\PYGZgt{}}\PYG{+w}{ }\PYG{p}{\PYGZob{}}\PYG{+w}{ }\PYG{n+nf+fm}{println!}\PYG{p}{(}\PYG{l+s}{\PYGZdq{}}\PYG{l+s}{integer 2}\PYG{l+s}{\PYGZdq{}}\PYG{p}{)}\PYG{p}{;}\PYG{+w}{ }\PYG{p}{\PYGZcb{}}
\PYG{+w}{        }\PYG{n}{\PYGZus{}}\PYG{+w}{ }\PYG{o}{=}\PYG{o}{\PYGZgt{}}\PYG{+w}{ }\PYG{p}{\PYGZob{}}\PYG{+w}{ }\PYG{n+nf+fm}{println!}\PYG{p}{(}\PYG{l+s}{\PYGZdq{}}\PYG{l+s}{something else}\PYG{l+s}{\PYGZdq{}}\PYG{p}{)}\PYG{p}{;}\PYG{+w}{ }\PYG{p}{\PYGZcb{}}
\PYG{+w}{    }\PYG{p}{\PYGZcb{}}
\PYG{p}{\PYGZcb{}}
\end{Verbatim}

与联合体变体匹配但不指定值的 \texttt{match} 分支永远都会成功。如果 \texttt{u} 的最后一个写入字段是 \texttt{u.i}，则以下代码将导致未定义行为：

\begin{Verbatim}[commandchars=\\\{\},breaklines=true,breakanywhere=true,fontsize=\codefont,xleftmargin=17bp]
\PYG{c+c1}{// 未定义行为！}
\PYG{k}{unsafe}\PYG{+w}{ }\PYG{p}{\PYGZob{}}
\PYG{+w}{    }\PYG{k}{match}\PYG{+w}{ }\PYG{n}{u}\PYG{+w}{ }\PYG{p}{\PYGZob{}}
\PYG{+w}{        }\PYG{n}{FloatOrInt}\PYG{+w}{ }\PYG{p}{\PYGZob{}}\PYG{+w}{ }\PYG{n}{f}\PYG{+w}{ }\PYG{p}{\PYGZcb{}}\PYG{+w}{ }\PYG{o}{=}\PYG{o}{\PYGZgt{}}\PYG{+w}{ }\PYG{p}{\PYGZob{}}\PYG{+w}{ }\PYG{n+nf+fm}{println!}\PYG{p}{(}\PYG{l+s}{\PYGZdq{}}\PYG{l+s}{float \PYGZob{}\PYGZcb{}}\PYG{l+s}{\PYGZdq{}}\PYG{p}{,}\PYG{+w}{ }\PYG{n}{f}\PYG{p}{)}\PYG{+w}{ }\PYG{p}{\PYGZcb{}}\PYG{p}{,}
\PYG{+w}{        }\PYG{c+c1}{// 警告：无法抵达的模式}
\PYG{+w}{        }\PYG{n}{FloatOrInt}\PYG{+w}{ }\PYG{p}{\PYGZob{}}\PYG{+w}{ }\PYG{n}{i}\PYG{+w}{ }\PYG{p}{\PYGZcb{}}\PYG{+w}{ }\PYG{o}{=}\PYG{o}{\PYGZgt{}}\PYG{+w}{ }\PYG{p}{\PYGZob{}}\PYG{+w}{ }\PYG{n+nf+fm}{println!}\PYG{p}{(}\PYG{l+s}{\PYGZdq{}}\PYG{l+s}{int \PYGZob{}\PYGZcb{}}\PYG{l+s}{\PYGZdq{}}\PYG{p}{,}\PYG{+w}{ }\PYG{n}{i}\PYG{p}{)}\PYG{+w}{ }\PYG{p}{\PYGZcb{}}
\PYG{+w}{    }\PYG{p}{\PYGZcb{}}
\PYG{p}{\PYGZcb{}}
\end{Verbatim}

\section{借用联合体}

借用联合体的一个字段就是借用整个联合体。这意味着，按照正常的借用规则，将一个字段作为可变借用会排斥对该字段或其他字段的任何其他借用，而将一个字段作为不可变借用则意味着对任何字段都不能再进行可变借用。

正如我们将在第 23 章中看到的，Rust 不仅可以帮你为自己的不安全代码构建出安全接口，还可以为用其他语言编写的代码构建出安全接口。从字面来看，“不安全”是充满危险的，但如果谨慎使用，那么也可以构建出高性能代码，同时还能让 Rust 程序员继续享有安全感。

